\documentclass[8pt]{beamer}
\usetheme{Madrid}
\usecolortheme{default}
\usepackage{amsmath, amssymb, graphicx, array, multirow, booktabs, xcolor, tikz}
\usetikzlibrary{shapes,arrows,positioning,fit,backgrounds}

% Compact font settings
\setbeamerfont{title}{size=\Large}
\setbeamerfont{subtitle}{size=\small}
\setbeamerfont{author}{size=\scriptsize}
\setbeamerfont{date}{size=\scriptsize}
\setbeamerfont{frametitle}{size=\large}
\setbeamerfont{framesubtitle}{size=\normalsize}
\setbeamerfont{enumerate item}{size=\footnotesize}
\setbeamerfont{itemize item}{size=\footnotesize}
\setbeamerfont{block title}{size=\footnotesize}

% Compact spacing
\setlength{\itemsep}{0.08ex}
\setlength{\parskip}{0.08ex}
\setlength{\leftmargini}{0.3cm}
\setbeamersize{text margin left=3mm, text margin right=3mm}
\addtobeamertemplate{frame begin}{\vspace{-0.4cm}}{}

% Colors
\definecolor{myblue}{RGB}{0,0,255}
\definecolor{myred}{RGB}{255,0,0}
\definecolor{mygreen}{RGB}{0,128,0}
\definecolor{myorange}{RGB}{255,165,0}
\definecolor{mypurple}{RGB}{128,0,128}
\definecolor{mygray}{RGB}{128,128,128}
\definecolor{mygold}{RGB}{255,215,0}

\title{Supervised Learning - Model Estimation}
\subtitle{100 Questions: OLS, MLE, Gradient Descent, Regularization, Cross-Validation}
\author{GARP RAI Exam Preparation}
\date{}

\begin{document}
	
	\begin{frame}
		\titlepage
	\end{frame}
	
	% ============================================================
	% SECTION 1: ESTIMATION PARADIGMS - QUESTIONS 1-15
	% ============================================================
	
	\begin{frame}{Q1: OLS Definition - Tutorial}
		\fontsize{6.5}{7.5}\selectfont
		\textbf{What is Ordinary Least Squares (OLS) estimation in supervised learning?}
		
		\vspace{0.05cm}
		\textcolor{myblue}{\textbf{ANSWER: B}} - A method that minimizes the sum of squared errors between predicted and actual values.
		
		\vspace{0.05cm}
		\textbf{DETAILED EXPLANATION:}
		\begin{itemize}
			\item OLS is the foundational estimation method in regression analysis
			\item Objective function: $\min \sum_{i=1}^N (y_i - \hat{y}_i)^2 = \min \sum_{i=1}^N \varepsilon_i^2$
			\item $\hat{y}_i = \beta_0 + \beta_1 x_{i1} + ... + \beta_p x_{ip}$ is the predicted value
			\item $\varepsilon_i = y_i - \hat{y}_i$ is the residual (error)
			\item OLS minimizes the Residual Sum of Squares (RSS)
			\item Developed by Carl Friedrich Gauss in the 1790s
			\item Has a closed-form solution: $\hat{\beta} = (X^TX)^{-1}X^Ty$
		\end{itemize}
		
		\vspace{0.05cm}
		\textcolor{myred}{\textbf{WHY OTHER OPTIONS ARE WRONG:}}
		\begin{itemize}
			\item \textcolor{myred}{A:} This describes Least Absolute Deviations (LAD) - minimizes $\sum |\varepsilon_i|$ instead of $\sum \varepsilon_i^2$
			\item \textcolor{myred}{C:} This describes Maximum Likelihood Estimation (MLE) - different paradigm entirely
			\item \textcolor{myred}{D:} This describes regularization (L1/Lasso) - penalizes complexity
		\end{itemize}
		
		\textcolor{myred}{\textbf{EXAM TRAP:}} OLS minimizes \textcolor{myblue}{squared errors}, NOT absolute errors! This is the most common confusion.
	\end{frame}
	
	\begin{frame}{Q2: OLS Objective Function - Tutorial}
		\fontsize{6.5}{7.5}\selectfont
		\textbf{Which is the correct objective function for OLS?}
		
		\vspace{0.05cm}
		\textcolor{myblue}{\textbf{ANSWER: B}} - $L = \sum_{i=1}^N (y_i - \hat{y}_i)^2$
		
		\vspace{0.05cm}
		\textbf{DETAILED EXPLANATION:}
		\begin{itemize}
			\item The OLS objective function is the Sum of Squared Errors (SSE)
			\item Also called Residual Sum of Squares (RSS) or Sum of Squared Residuals (SSR)
			\item Mathematical notation: $L(\beta) = \sum_{i=1}^N (y_i - X_i\beta)^2$
			\item This function is convex and differentiable $\Rightarrow$ has unique minimum
			\item First derivative: $\frac{\partial L}{\partial \beta} = -2X^T(y - X\beta)$
			\item Setting derivative to zero yields normal equations: $X^TX\beta = X^Ty$
		\end{itemize}
		
		\vspace{0.05cm}
		\textcolor{myred}{\textbf{WHY OTHER OPTIONS ARE WRONG:}}
		\begin{itemize}
			\item \textcolor{myred}{A:} L1 loss (Least Absolute Deviations) - not differentiable at zero
			\item \textcolor{myred}{C:} L4 loss - rarely used, not OLS
			\item \textcolor{myred}{D:} Chebyshev loss (minimax) - minimizes maximum error
		\end{itemize}
		
		\textcolor{myred}{\textbf{EXAM TRAP:}} OLS = \textcolor{myblue}{RSS} = $\sum (y_i - \hat{y}_i)^2$! Know all three names (SSE, RSS, SSR).
	\end{frame}
	
	\begin{frame}{Q3: OLS vs NLS - Tutorial}
		\fontsize{6.5}{7.5}\selectfont
		\textbf{Key difference between OLS and Nonlinear Least Squares?}
		
		\vspace{0.05cm}
		\textcolor{myblue}{\textbf{ANSWER: B}} - OLS is used when model is linear in parameters; NLS when nonlinear in parameters.
		
		\vspace{0.05cm}
		\textbf{DETAILED EXPLANATION:}
		\begin{itemize}
			\item \textbf{Linear in parameters:} Model can be written as $\beta_0 + \beta_1 f_1(x) + \beta_2 f_2(x) + ...$
			\begin{itemize}
				\item Example: $y = \beta_0 + \beta_1 x + \beta_2 x^2$ (polynomial regression)
				\item Even though nonlinear in x, it's linear in $\beta$s
			\end{itemize}
			\item \textbf{Nonlinear in parameters:} Parameters appear nonlinearly
			\begin{itemize}
				\item Example: $y = \beta_0 + \beta_1 e^{\beta_2 x}$
				\item $\beta_2$ appears in exponent - nonlinear
			\end{itemize}
			\item OLS has closed-form solution: $\hat{\beta} = (X^TX)^{-1}X^Ty$
			\item NLS requires iterative numerical optimization
		\end{itemize}
		
		\vspace{0.05cm}
		\textcolor{myred}{\textbf{WHY OTHER OPTIONS ARE WRONG:}}
		\begin{itemize}
			\item \textcolor{myred}{A:} Both minimize squared errors - this is the similarity, not difference
			\item \textcolor{myred}{C:} Actually reversed - OLS has closed-form, NLS needs iterative methods
			\item \textcolor{myred}{D:} Both are regression methods (continuous dependent variables)
		\end{itemize}
		
		\textcolor{myred}{\textbf{EXAM TRAP:}} Difference is \textcolor{myblue}{linearity in parameters}, NOT in variables! $y = \beta_0 + \beta_1 x^2$ is OLS.
	\end{frame}
	
	\begin{frame}{Q4: NLS Example - Tutorial}
		\fontsize{6.5}{7.5}\selectfont
		\textbf{Which model requires Nonlinear Least Squares (NLS) estimation?}
		
		\vspace{0.05cm}
		\textcolor{myblue}{\textbf{ANSWER: C}} - $y = \beta_0 + \beta_1 e^{\beta_2 x} + \varepsilon$
		
		\vspace{0.05cm}
		\textbf{DETAILED EXPLANATION:}
		\begin{itemize}
			\item Model C has $\beta_2$ in the exponent: $e^{\beta_2 x}$
			\item This is nonlinear in parameter $\beta_2$ - cannot be transformed to linear form
			\item Parameter $\beta_1$ multiplies the exponential term
			\item NLS requires iterative methods like Gauss-Newton or Levenberg-Marquardt
			\item The normal equations become nonlinear and cannot be solved directly
		\end{itemize}
		
		\vspace{0.05cm}
		\textcolor{myred}{\textbf{WHY OTHER OPTIONS ARE WRONG:}}
		\begin{itemize}
			\item \textcolor{myred}{A:} $y = \beta_0 + \beta_1 x_1 + \beta_2 x_2$ - linear in all parameters (OLS)
			\item \textcolor{myred}{B:} $y = \beta_0 + \beta_1 x + \beta_2 x^2$ - linear in $\beta$s (polynomial regression, OLS)
			\item \textcolor{myred}{D:} $y = \beta_0 + \beta_1 \ln(x)$ - linear in $\beta$s (OLS with transformed x)
		\end{itemize}
		
		\textcolor{myred}{\textbf{EXAM TRAP:}} NLS is about \textcolor{myblue}{nonlinearity in parameters}, not variables! This is the most common mistake.
	\end{frame}
	
	\begin{frame}{Q5: MLE Definition - Tutorial}
		\fontsize{6.5}{7.5}\selectfont
		\textbf{What is Maximum Likelihood Estimation (MLE)?}
		
		\vspace{0.05cm}
		\textcolor{myblue}{\textbf{ANSWER: B}} - Method that finds parameters maximizing probability (likelihood) of observing given data.
		
		\vspace{0.05cm}
		\textbf{DETAILED EXPLANATION:}
		\begin{itemize}
			\item MLE finds $\theta$ that maximizes $L(\theta) = P(\text{data} | \theta)$
			\item \textbf{Intuition:} "What parameters make the observed data most probable?"
			\item Likelihood function: $L(\theta) = \prod_{i=1}^N f(y_i | \theta)$ for i.i.d. data
			\item Log-likelihood: $\ell(\theta) = \sum_{i=1}^N \log f(y_i | \theta)$
			\item MLE is asymptotically consistent, efficient, and normally distributed
			\item Under regularity conditions, MLE achieves Cramér-Rao lower bound asymptotically
		\end{itemize}
		
		\vspace{0.05cm}
		\textcolor{myred}{\textbf{WHY OTHER OPTIONS ARE WRONG:}}
		\begin{itemize}
			\item \textcolor{myred}{A:} This describes OLS - different estimation paradigm
			\item \textcolor{myred}{C:} This describes regularization (penalizing complexity)
			\item \textcolor{myred}{D:} Minimizing error variance is related to efficiency but not MLE definition
		\end{itemize}
		
		\textcolor{myred}{\textbf{EXAM TRAP:}} MLE maximizes \textcolor{myblue}{likelihood} $P(\text{data}|\theta)$, NOT $P(\theta|\text{data})$!
	\end{frame}
	
	\begin{frame}{Q6: OLS vs MLE Equivalence - Tutorial}
		\fontsize{6.5}{7.5}\selectfont
		\textbf{When do OLS and MLE produce same parameter estimates?}
		
		\vspace{0.05cm}
		\textcolor{myblue}{\textbf{ANSWER: B}} - When errors are normally distributed with constant variance.
		
		\vspace{0.05cm}
		\textbf{DETAILED EXPLANATION:}
		\begin{itemize}
			\item Model: $y_i = X_i\beta + \varepsilon_i$, with $\varepsilon_i \sim N(0, \sigma^2)$
			\item Likelihood: $L(\beta, \sigma^2) = \prod_{i=1}^N \frac{1}{\sqrt{2\pi\sigma^2}} \exp\left(-\frac{(y_i - X_i\beta)^2}{2\sigma^2}\right)$
			\item Log-likelihood: $\ell(\beta) = -\frac{N}{2}\log(2\pi) - \frac{N}{2}\log(\sigma^2) - \frac{1}{2\sigma^2}\sum_{i=1}^N (y_i - X_i\beta)^2$
			\item Maximizing $\ell(\beta)$ is equivalent to minimizing $\sum (y_i - X_i\beta)^2$ (RSS)
			\item The $\sigma^2$ term doesn't affect the $\beta$ estimates
		\end{itemize}
		
		\vspace{0.05cm}
		\textcolor{myred}{\textbf{WHY OTHER OPTIONS ARE WRONG:}}
		\begin{itemize}
			\item \textcolor{myred}{A:} Independence alone insufficient - need normal distribution
			\item \textcolor{myred}{C:} Uniform errors would give different estimator (MLE would minimize maximum absolute error)
			\item \textcolor{myred}{D:} Autocorrelation violates OLS assumptions
		\end{itemize}
		
		\textcolor{myred}{\textbf{EXAM TRAP:}} OLS and MLE are \textcolor{myblue}{equivalent} under normal errors! This is a critical exam concept.
	\end{frame}
	
	\begin{frame}{Q7: Logistic Regression Estimation - Tutorial}
		\fontsize{6.5}{7.5}\selectfont
		\textbf{What estimation method is typically used for logistic regression?}
		
		\vspace{0.05cm}
		\textcolor{myblue}{\textbf{ANSWER: C}} - Maximum Likelihood Estimation (MLE)
		
		\vspace{0.05cm}
		\textbf{DETAILED EXPLANATION:}
		\begin{itemize}
			\item Logistic regression: $\log\left(\frac{p_i}{1-p_i}\right) = X_i\beta$
			\item $p_i = P(Y_i = 1 | X_i) = \frac{1}{1 + e^{-X_i\beta}}$
			\item Target is binary: $Y_i \in \{0, 1\}$
			\item Bernoulli likelihood: $L(\beta) = \prod_{i=1}^N p_i^{Y_i}(1-p_i)^{1-Y_i}$
			\item Log-likelihood: $\ell(\beta) = \sum_{i=1}^N [Y_i\log(p_i) + (1-Y_i)\log(1-p_i)]$
			\item No closed-form solution - requires numerical optimization (Newton-Raphson, IRLS)
		\end{itemize}
		
		\vspace{0.05cm}
		\textcolor{myred}{\textbf{WHY OTHER OPTIONS ARE WRONG:}}
		\begin{itemize}
			\item \textcolor{myred}{A:} OLS inappropriate - binary target, non-normal errors, heteroscedasticity
			\item \textcolor{myred}{B:} NLS not standard - logistic regression's nonlinearity handled by MLE
			\item \textcolor{myred}{D:} Method of Moments less common and less efficient
		\end{itemize}
		
		\textcolor{myred}{\textbf{EXAM TRAP:}} Logistic regression uses \textcolor{myblue}{MLE}, NOT OLS! Critical distinction.
	\end{frame}
	
	\begin{frame}{Q8: Likelihood Function - Tutorial}
		\fontsize{6.5}{7.5}\selectfont
		\textbf{What does $L(\theta | \text{data})$ represent?}
		
		\vspace{0.05cm}
		\textcolor{myblue}{\textbf{ANSWER: B}} - The probability of the data given the parameters.
		
		\vspace{0.05cm}
		\textbf{DETAILED EXPLANATION:}
		\begin{itemize}
			\item $L(\theta | \text{data}) = P(\text{data} | \theta)$
			\item \textbf{Key distinction:} Likelihood vs Posterior
			\begin{itemize}
				\item Likelihood: $P(\text{data} | \theta)$ - frequentist paradigm
				\item Posterior: $P(\theta | \text{data})$ - Bayesian paradigm
			\end{itemize}
			\item Likelihood treats $\theta$ as fixed unknown constant
			\item Data is random; parameters are fixed
			\item Uses Bayes' theorem: $P(\theta | \text{data}) \propto P(\text{data} | \theta) P(\theta)$
		\end{itemize}
		
		\vspace{0.05cm}
		\textcolor{myred}{\textbf{WHY OTHER OPTIONS ARE WRONG:}}
		\begin{itemize}
			\item \textcolor{myred}{A:} This describes posterior distribution $P(\theta|\text{data})$ - Bayesian concept
			\item \textcolor{myred}{C:} This is the OLS loss function (RSS) - different concept
			\item \textcolor{myred}{D:} This describes estimator variance - different statistical concept
		\end{itemize}
		
		\textcolor{myred}{\textbf{EXAM TRAP:}} Likelihood = $P(\text{data} | \theta)$, NOT $P(\theta | \text{data})$! This confusion is extremely common.
	\end{frame}
	
	\begin{frame}{Q9: Log-Likelihood - Tutorial}
		\fontsize{6.5}{7.5}\selectfont
		\textbf{Why is log-likelihood often used instead of likelihood?}
		
		\vspace{0.05cm}
		\textcolor{myblue}{\textbf{ANSWER: B}} - Easier to maximize (products $\to$ sums, monotonic).
		
		\vspace{0.05cm}
		\textbf{DETAILED EXPLANATION:}
		\begin{itemize}
			\item \textbf{Mathematical convenience:}
			\begin{itemize}
				\item Likelihood involves products: $L(\theta) = \prod_{i=1}^N f(y_i|\theta)$
				\item Log-likelihood involves sums: $\ell(\theta) = \sum_{i=1}^N \log f(y_i|\theta)$
			\end{itemize}
			\item \textbf{Monotonicity:} $\log$ is strictly increasing
			\begin{itemize}
				\item Maximizing $\ell(\theta)$ $\equiv$ maximizing $L(\theta)$
				\item If $L(\theta_1) > L(\theta_2)$, then $\ell(\theta_1) > \ell(\theta_2)$
			\end{itemize}
			\item \textbf{Numerical stability:}
			\begin{itemize}
				\item Products of small probabilities can underflow to zero
				\item Sums of logs are more numerically stable
			\end{itemize}
			\item \textbf{Derivative simplification:} Sums easier to differentiate than products
		\end{itemize}
		
		\vspace{0.05cm}
		\textcolor{myred}{\textbf{WHY OTHER OPTIONS ARE WRONG:}}
		\begin{itemize}
			\item \textcolor{myred}{A:} False - log-likelihood is often negative (probabilities < 1)
			\item \textcolor{myred}{C:} False - usually requires numerical optimization
			\item \textcolor{myred}{D:} False - doesn't affect parameter count (that's regularization)
		\end{itemize}
		
		\textcolor{myred}{\textbf{EXAM TRAP:}} Log-likelihood is \textcolor{myblue}{monotonic} - maximizing it $\equiv$ maximizing likelihood!
	\end{frame}
	
	\begin{frame}{Q10: Closed-Form vs Numerical - Tutorial}
		\fontsize{6.5}{7.5}\selectfont
		\textbf{Which estimation method typically has a closed-form solution?}
		
		\vspace{0.05cm}
		\textcolor{myblue}{\textbf{ANSWER: C}} - Ordinary Least Squares for linear regression.
		
		\vspace{0.05cm}
		\textbf{DETAILED EXPLANATION:}
		\begin{itemize}
			\item \textbf{OLS closed-form:} $\hat{\beta} = (X^TX)^{-1}X^Ty$
			\begin{itemize}
				\item Derived from normal equations: $X^TX\hat{\beta} = X^Ty$
				\item Requires $X^TX$ to be invertible (no perfect multicollinearity)
				\item Computational complexity: $O(p^2N + p^3)$
			\end{itemize}
			\item \textbf{Why others don't have closed-form:}
			\begin{itemize}
				\item Logistic regression: nonlinear in parameters (link function)
				\item NLS: nonlinear in parameters
				\item Neural networks: highly complex composition of functions
			\end{itemize}
			\item \textbf{When closed-form is not possible:}
			\begin{itemize}
				\item When normal equations are nonlinear
				\item When likelihood has no analytical solution
				\item When loss function is not convex
			\end{itemize}
		\end{itemize}
		
		\vspace{0.05cm}
		\textcolor{myred}{\textbf{WHY OTHER OPTIONS ARE WRONG:}}
		\begin{itemize}
			\item \textcolor{myred}{A:} Logistic regression MLE requires numerical optimization (Newton-Raphson)
			\item \textcolor{myred}{B:} NLS requires iterative numerical optimization
			\item \textcolor{myred}{D:} Neural networks use backpropagation (numerical)
		\end{itemize}
		
		\textcolor{myred}{\textbf{EXAM TRAP:}} Only \textcolor{myblue}{OLS for linear regression} has a closed-form solution!
	\end{frame}
	
	\begin{frame}{Q11: OLS Assumptions - Tutorial}
		\fontsize{6.5}{7.5}\selectfont
		\textbf{Which is NOT a standard OLS assumption?}
		
		\vspace{0.05cm}
		\textcolor{myblue}{\textbf{ANSWER: D}} - The errors are autocorrelated.
		
		\vspace{0.05cm}
		\textbf{DETAILED EXPLANATION:}
		\begin{itemize}
			\item \textbf{Standard OLS assumptions (Gauss-Markov):}
			\begin{enumerate}
				\item \textbf{Linearity in parameters:} $y = X\beta + \varepsilon$
				\item \textbf{Zero conditional mean:} $E[\varepsilon | X] = 0$
				\item \textbf{No perfect multicollinearity:} $X^TX$ invertible
				\item \textbf{Homoscedasticity:} $Var(\varepsilon | X) = \sigma^2$ (constant)
				\item \textbf{No autocorrelation:} $Cov(\varepsilon_i, \varepsilon_j | X) = 0$ for $i \neq j$
				\item \textbf{Normality:} $\varepsilon \sim N(0, \sigma^2)$ (for inference)
			\end{enumerate}
			\item \textbf{Autocorrelation} is a violation of the independence assumption
			\begin{itemize}
				\item Common in time series data
				\item Leads to inefficient estimates (not BLUE)
				\item Standard errors are biased
			\end{itemize}
		\end{itemize}
		
		\vspace{0.05cm}
		\textcolor{myred}{\textbf{WHY OTHER OPTIONS ARE WRONG:}}
		\begin{itemize}
			\item \textcolor{myred}{A:} This IS an OLS assumption
			\item \textcolor{myred}{B:} This IS an OLS assumption
			\item \textcolor{myred}{C:} This IS an OLS assumption (for inference)
		\end{itemize}
		
		\textcolor{myred}{\textbf{EXAM TRAP:}} OLS assumes \textcolor{myblue}{no autocorrelation} in errors! Autocorrelation is a violation.
	\end{frame}
	
	\begin{frame}{Q12: MLE Properties - Tutorial}
		\fontsize{6.5}{7.5}\selectfont
		\textbf{Which is a desirable property of Maximum Likelihood Estimators?}
		
		\vspace{0.05cm}
		\textcolor{myblue}{\textbf{ANSWER: B}} - They are consistent (converge to true values as sample size increases).
		
		\vspace{0.05cm}
		\textbf{DETAILED EXPLANATION:}
		\begin{itemize}
			\item \textbf{Consistency:} $\hat{\theta}_{MLE} \xrightarrow{p} \theta_0$ as $n \to \infty$
			\begin{itemize}
				\item Under regularity conditions (identifiability, differentiability, etc.)
				\item The estimator converges in probability to the true parameter
			\end{itemize}
			\item \textbf{Asymptotic Normality:} $\sqrt{n}(\hat{\theta}_{MLE} - \theta_0) \xrightarrow{d} N(0, I^{-1}(\theta_0))$
			\begin{itemize}
				\item $I(\theta)$ is the Fisher Information matrix
				\item Achieves Cramér-Rao lower bound asymptotically
			\end{itemize}
			\item \textbf{Asymptotic Efficiency:} Minimum variance among regular estimators
			\item \textbf{Invariance:} $\widehat{g(\theta)} = g(\hat{\theta})$ for continuous $g$
		\end{itemize}
		
		\vspace{0.05cm}
		\textcolor{myred}{\textbf{WHY OTHER OPTIONS ARE WRONG:}}
		\begin{itemize}
			\item \textcolor{myred}{A:} False - MLE can be biased (e.g., $\hat{\sigma}^2_{MLE}$ is biased for $\sigma^2$)
			\item \textcolor{myred}{C:} False - asymptotically efficient, not always minimum variance in finite samples
			\item \textcolor{myred}{D:} False - often requires numerical optimization
		\end{itemize}
		
		\textcolor{myred}{\textbf{EXAM TRAP:}} MLE is \textcolor{myblue}{consistent} but NOT necessarily unbiased!
	\end{frame}
	
	\begin{frame}{Q13: NLS Optimization - Tutorial}
		\fontsize{6.5}{7.5}\selectfont
		\textbf{Why does NLS require iterative numerical optimization?}
		
		\vspace{0.05cm}
		\textcolor{myblue}{\textbf{ANSWER: B}} - No closed-form solution for nonlinear-in-parameters models.
		
		\vspace{0.05cm}
		\textbf{DETAILED EXPLANATION:}
		\begin{itemize}
			\item \textbf{The problem:} Normal equations are nonlinear
			\begin{itemize}
				\item OLS: $X^TX\beta = X^Ty$ $\Rightarrow$ linear system
				\item NLS: $\frac{\partial S(\beta)}{\partial \beta} = 0$ where $S(\beta) = \sum (y_i - f(x_i, \beta))^2$
				\item Partial derivatives $\frac{\partial S}{\partial \beta_j}$ are nonlinear functions of $\beta$
			\end{itemize}
			\item \textbf{Common NLS algorithms:}
			\begin{itemize}
				\item Gauss-Newton: Linearizes the model around current estimates
				\item Levenberg-Marquardt: Combines Gauss-Newton and gradient descent
				\item Gradient descent (first-order method)
				\item Newton-Raphson (second-order method)
			\end{itemize}
			\item \textbf{Convergence issues:}
			\begin{itemize}
				\item Requires good starting values
				\item May converge to local minima
				\item May not converge at all if model is misspecified
			\end{itemize}
		\end{itemize}
		
		\vspace{0.05cm}
		\textcolor{myred}{\textbf{WHY OTHER OPTIONS ARE WRONG:}}
		\begin{itemize}
			\item \textcolor{myred}{A:} Multiple local minima are a concern but not the primary reason
			\item \textcolor{myred}{C:} Error distribution doesn't determine need for numerical optimization
			\item \textcolor{myred}{D:} Parameter count affects complexity but isn't fundamental
		\end{itemize}
		
		\textcolor{myred}{\textbf{EXAM TRAP:}} No closed-form solution exists for \textcolor{myblue}{nonlinear-in-parameters} models!
	\end{frame}
	
	\begin{frame}{Q14: Gauss-Markov Theorem - Tutorial}
		\fontsize{6.5}{7.5}\selectfont
		\textbf{What does the Gauss-Markov theorem state about OLS estimators?}
		
		\vspace{0.05cm}
		\textcolor{myblue}{\textbf{ANSWER: B}} - OLS estimators are the Best Linear Unbiased Estimators (BLUE) under certain conditions.
		
		\vspace{0.05cm}
		\textbf{DETAILED EXPLANATION:}
		\begin{itemize}
			\item \textbf{Gauss-Markov assumptions:}
			\begin{enumerate}
				\item Linearity in parameters
				\item No perfect multicollinearity
				\item Zero conditional mean: $E[\varepsilon | X] = 0$
				\item Homoscedasticity: $Var(\varepsilon | X) = \sigma^2$
				\item No autocorrelation: $Cov(\varepsilon_i, \varepsilon_j | X) = 0$
			\end{enumerate}
			\item \textbf{BLUE means:}
			\begin{itemize}
				\item \textbf{Best:} Minimum variance among all linear unbiased estimators
				\item \textbf{Linear:} Estimator is a linear function of $y$
				\item \textbf{Unbiased:} $E[\hat{\beta}] = \beta$
				\item \textbf{Estimator:} $\hat{\beta} = (X^TX)^{-1}X^Ty$
			\end{itemize}
			\item \textbf{Important:} OLS is BLUE among \textit{linear} unbiased estimators
			\begin{itemize}
				\item Nonlinear estimators could have lower variance
				\item Biased estimators could have lower MSE (bias-variance tradeoff)
			\end{itemize}
		\end{itemize}
		
		\vspace{0.05cm}
		\textcolor{myred}{\textbf{WHY OTHER OPTIONS ARE WRONG:}}
		\begin{itemize}
			\item \textcolor{myred}{A:} Close but incomplete - misses "Best" (minimum variance)
			\item \textcolor{myred}{C:} False - only among linear unbiased estimators
			\item \textcolor{myred}{D:} False - consistency requires additional conditions
		\end{itemize}
		
		\textcolor{myred}{\textbf{EXAM TRAP:}} OLS is \textcolor{myblue}{BLUE} under Gauss-Markov assumptions! Not "always best" among all estimators.
	\end{frame}
	
	\begin{frame}{Q15: Estimation Methods Summary - Tutorial}
		\fontsize{6.5}{7.5}\selectfont
		\textbf{Which estimation method for $y = \beta_0 + \beta_1 x_1^{\beta_2} + \varepsilon$?}
		
		\vspace{0.05cm}
		\textcolor{myblue}{\textbf{ANSWER: D}} - Both NLS and MLE would be appropriate.
		
		\vspace{0.05cm}
		\textbf{DETAILED EXPLANATION:}
		\begin{itemize}
			\item \textbf{Model analysis:}
			\begin{itemize}
				\item $y = \beta_0 + \beta_1 x_1^{\beta_2} + \varepsilon$
				\item $\beta_2$ appears in the exponent of $x_1$
				\item This is nonlinear in parameter $\beta_2$
				\item Cannot be transformed to linear form (unlike $x^2$ which is linear in parameters)
			\end{itemize}
			\item \textbf{NLS approach:}
			\begin{itemize}
				\item Minimizes $\sum (y_i - \beta_0 - \beta_1 x_{1i}^{\beta_2})^2$
				\item Requires numerical optimization
				\item Equivalent to MLE if errors are normal
			\end{itemize}
			\item \textbf{MLE approach:}
			\begin{itemize}
				\item Defines likelihood: $L(\beta_0, \beta_1, \beta_2, \sigma^2)$
				\item If $\varepsilon \sim N(0, \sigma^2)$, MLE $\equiv$ NLS
				\item If errors not normal, MLE differs from NLS
			\end{itemize}
		\end{itemize}
		
		\vspace{0.05cm}
		\textcolor{myred}{\textbf{WHY OTHER OPTIONS ARE WRONG:}}
		\begin{itemize}
			\item \textcolor{myred}{A:} OLS requires linearity in parameters - this model violates that
			\item \textcolor{myred}{B:} NLS alone correct but incomplete - MLE also works
			\item \textcolor{myred}{C:} MLE alone correct but incomplete - NLS also works
		\end{itemize}
		
		\textcolor{myred}{\textbf{EXAM TRAP:}} For nonlinear models, both \textcolor{myblue}{NLS and MLE} can be used!
	\end{frame}
	
	% ============================================================
	% SECTION 2: GRADIENT DESCENT - QUESTIONS 16-35
	% ============================================================
	
	\begin{frame}{Q16: Gradient Descent Definition - Tutorial}
		\fontsize{6.5}{7.5}\selectfont
		\textbf{What is gradient descent in model estimation?}
		
		\vspace{0.05cm}
		\textcolor{myblue}{\textbf{ANSWER: B}} - Iterative optimization algorithm updating parameters in direction of negative gradient.
		
		\vspace{0.05cm}
		\textbf{DETAILED EXPLANATION:}
		\begin{itemize}
			\item \textbf{Core concept:} Iteratively move in the direction of steepest descent
			\item \textbf{Update rule:} $\theta_{new} = \theta_{old} - \alpha \nabla L(\theta)$
			\item \textbf{Components:}
			\begin{itemize}
				\item $\theta$: Parameters being optimized
				\item $\alpha$: Learning rate (step size)
				\item $\nabla L(\theta)$: Gradient (vector of partial derivatives)
			\end{itemize}
			\item \textbf{Intuition:}
			\begin{itemize}
				\item Gradient points in direction of steepest increase
				\item Negative gradient points in direction of steepest decrease
				\item Move in negative gradient direction to reduce loss
			\end{itemize}
			\item \textbf{Convergence conditions:}
			\begin{itemize}
				\item Loss function is differentiable
				\item Learning rate is chosen appropriately
				\item Function is convex for guaranteed global convergence
			\end{itemize}
		\end{itemize}
		
		\vspace{0.05cm}
		\textcolor{myred}{\textbf{WHY OTHER OPTIONS ARE WRONG:}}
		\begin{itemize}
			\item \textcolor{myred}{A:} False - can get stuck in local minima
			\item \textcolor{myred}{C:} This describes closed-form solutions
			\item \textcolor{myred}{D:} This describes random search
		\end{itemize}
		
		\textcolor{myred}{\textbf{EXAM TRAP:}} Gradient descent is \textcolor{myblue}{iterative} and uses \textcolor{myblue}{negative gradient}!
	\end{frame}
	
	\begin{frame}{Q17: Gradient Descent Update Rule - Tutorial}
		\fontsize{6.5}{7.5}\selectfont
		\textbf{What is the general update rule for gradient descent?}
		
		\vspace{0.05cm}
		\textcolor{myblue}{\textbf{ANSWER: B}} - $\theta_{new} = \theta_{old} - \alpha \nabla L(\theta)$
		
		\vspace{0.05cm}
		\textbf{DETAILED EXPLANATION:}
		\begin{itemize}
			\item \textbf{Formula breakdown:}
			\begin{itemize}
				\item $\theta_{new}$: Updated parameter values
				\item $\theta_{old}$: Current parameter values
				\item $\alpha$: Learning rate (hyperparameter)
				\item $\nabla L(\theta)$: Gradient of loss function
			\end{itemize}
			\item \textbf{Why subtract?}
			\begin{itemize}
				\item $\nabla L(\theta)$ points uphill (direction of increase)
				\item $-\nabla L(\theta)$ points downhill (direction of decrease)
				\item Subtracting moves toward lower loss
			\end{itemize}
			\item \textbf{Example for linear regression:}
			\begin{itemize}
				\item $L(\beta) = \sum (y_i - X_i\beta)^2$
				\item $\nabla L(\beta) = -2X^T(y - X\beta)$
				\item $\beta_{new} = \beta_{old} + 2\alpha X^T(y - X\beta)$
			\end{itemize}
		\end{itemize}
		
		\vspace{0.05cm}
		\textcolor{myred}{\textbf{WHY OTHER OPTIONS ARE WRONG:}}
		\begin{itemize}
			\item \textcolor{myred}{A:} This is gradient ascent (moving to increase loss)
			\item \textcolor{myred}{C:} Multiplication is not the correct operation
			\item \textcolor{myred}{D:} Division is not the correct operation
		\end{itemize}
		
		\textcolor{myred}{\textbf{EXAM TRAP:}} Gradient descent uses \textcolor{myblue}{subtraction} of the gradient!
	\end{frame}
	
	\begin{frame}{Q18: Learning Rate Definition - Tutorial}
		\fontsize{6.5}{7.5}\selectfont
		\textbf{What does the learning rate $\alpha$ control in gradient descent?}
		
		\vspace{0.05cm}
		\textcolor{myblue}{\textbf{ANSWER: B}} - The step size of each parameter update.
		
		\vspace{0.05cm}
		\textbf{DETAILED EXPLANATION:}
		\begin{itemize}
			\item \textbf{Step size control:}
			\begin{itemize}
				\item $\alpha$ determines how far we move in the gradient direction
				\item Larger $\alpha$: larger steps, faster potential convergence
				\item Smaller $\alpha$: smaller steps, slower but more stable
			\end{itemize}
			\item \textbf{Effects of learning rate:}
			\begin{itemize}
				\item \textbf{Too large:} Oscillation, divergence, failure to converge
				\item \textbf{Too small:} Very slow convergence, getting stuck in plateaus
				\item \textbf{Just right:} Efficient convergence to minimum
			\end{itemize}
			\item \textbf{Common learning rate values:}
			\begin{itemize}
				\item Typical range: $10^{-6}$ to $10^{-1}$
				\item Often use learning rate schedules (decay over time)
				\item Adaptive methods (Adam) adjust per parameter
			\end{itemize}
			\item \textbf{Learning rate vs other hyperparameters:}
			\begin{itemize}
				\item Different from number of epochs (iterations)
				\item Different from batch size
				\item Different from regularization parameter $\lambda$
			\end{itemize}
		\end{itemize}
		
		\vspace{0.05cm}
		\textcolor{myred}{\textbf{WHY OTHER OPTIONS ARE WRONG:}}
		\begin{itemize}
			\item \textcolor{myred}{A:} Direction determined by gradient sign, not $\alpha$
			\item \textcolor{myred}{C:} Number of iterations is a different hyperparameter
			\item \textcolor{myred}{D:} Initialization is a different step
		\end{itemize}
		
		\textcolor{myred}{\textbf{EXAM TRAP:}} Learning rate controls \textcolor{myblue}{step size}, NOT direction!
	\end{frame}
	
	\begin{frame}{Q19: Learning Rate Too Large - Tutorial}
		\fontsize{6.5}{7.5}\selectfont
		\textbf{What happens if learning rate is set too large?}
		
		\vspace{0.05cm}
		\textcolor{myblue}{\textbf{ANSWER: B}} - Algorithm may oscillate or diverge.
		
		\vspace{0.05cm}
		\textbf{DETAILED EXPLANATION:}
		\begin{itemize}
			\item \textbf{Consequences of large $\alpha$:}
			\begin{enumerate}
				\item \textbf{Oscillation:} Parameter jumps back and forth across minimum
				\item \textbf{Divergence:} Loss increases without bound
				\item \textbf{Instability:} Updates become erratic and unpredictable
			\end{enumerate}
			\item \textbf{Mathematical explanation:}
			\begin{itemize}
				\item Consider quadratic: $f(\theta) = \frac{1}{2}\theta^2$
				\item Gradient descent: $\theta_{t+1} = \theta_t - \alpha \theta_t = (1-\alpha)\theta_t$
				\item Convergence requires $|1-\alpha| < 1$ $\Rightarrow$ $0 < \alpha < 2$
				\item If $\alpha > 2$: divergence (oscillating with increasing magnitude)
				\item If $\alpha = 2$: oscillation between $\theta$ and $-\theta$
			\end{itemize}
			\item \textbf{Visual analogy:} Like walking down a hill with too large steps
			\begin{itemize}
				\item Overshoot the bottom
				\item End up higher than where you started
				\item Keep oscillating or fall off the mountain
			\end{itemize}
		\end{itemize}
		
		\vspace{0.05cm}
		\textcolor{myred}{\textbf{WHY OTHER OPTIONS ARE WRONG:}}
		\begin{itemize}
			\item \textcolor{myred}{A:} Too small learning rate causes slow convergence
			\item \textcolor{myred}{C:} Too large prevents convergence
			\item \textcolor{myred}{D:} Algorithm already deterministic with fixed initialization
		\end{itemize}
		
		\textcolor{myred}{\textbf{EXAM TRAP:}} Too large learning rate $\rightarrow$ \textcolor{myred}{divergence or oscillation}!
	\end{frame}
	
	\begin{frame}{Q20: Learning Rate Too Small - Tutorial}
		\fontsize{6.5}{7.5}\selectfont
		\textbf{What happens if learning rate is set too small?}
		
		\vspace{0.05cm}
		\textcolor{myblue}{\textbf{ANSWER: B}} - Algorithm converges very slowly and may get stuck.
		
		\vspace{0.05cm}
		\textbf{DETAILED EXPLANATION:}
		\begin{itemize}
			\item \textbf{Consequences of small $\alpha$:}
			\begin{enumerate}
				\item \textbf{Slow convergence:} Tiny steps require many iterations
				\item \textbf{Getting stuck:} May stop in plateaus or local minima
				\item \textbf{Wasted computation:} Many iterations with little progress
			\end{enumerate}
			\item \textbf{Mathematical explanation:}
			\begin{itemize}
				\item For convex function: $\theta_{t+1} = \theta_t - \alpha \nabla f(\theta_t)$
				\item If $\alpha$ is very small, $\theta_t$ changes very slowly
				\item Convergence rate $O(1/t)$ for gradient descent
				\item Small $\alpha$ increases the constant in $O(1/t)$
			\end{itemize}
			\item \textbf{Visual analogy:} Like walking down a hill with tiny steps
			\begin{itemize}
				\item Make very slow progress
				\item May appear to stop on flat ground (plateau)
				\item Will eventually reach bottom but takes forever
			\end{itemize}
			\item \textbf{Diagnosis:}
			\begin{itemize}
				\item Loss decreases very slowly
				\item Requires many more epochs than expected
				\item Often confused with convergence
			\end{itemize}
		\end{itemize}
		
		\vspace{0.05cm}
		\textcolor{myred}{\textbf{WHY OTHER OPTIONS ARE WRONG:}}
		\begin{itemize}
			\item \textcolor{myred}{A:} Too small gives slow convergence, not quick
			\item \textcolor{myred}{C:} No algorithm finds global minimum instantly
			\item \textcolor{myred}{D:} Too small is stable but slow
		\end{itemize}
		
		\textcolor{myred}{\textbf{EXAM TRAP:}} Too small learning rate $\rightarrow$ \textcolor{myred}{slow convergence and getting stuck}!
	\end{frame}
	
	\begin{frame}{Q21: Batch Gradient Descent - Tutorial}
		\fontsize{6.5}{7.5}\selectfont
		\textbf{What is Batch Gradient Descent?}
		
		\vspace{0.05cm}
		\textcolor{myblue}{\textbf{ANSWER: B}} - Uses entire training dataset per update.
		
		\vspace{0.05cm}
		\textbf{DETAILED EXPLANATION:}
		\begin{itemize}
			\item \textbf{Definition:} Computes gradient using all N training samples
			\begin{itemize}
				\item $\nabla L(\theta) = \frac{1}{N}\sum_{i=1}^N \nabla L_i(\theta)$
				\item One update per epoch (full pass through data)
			\end{itemize}
			\item \textbf{Characteristics:}
			\begin{itemize}
				\item \textbf{Accurate gradient:} Uses all data, no sampling noise
				\item \textbf{Stable:} Smooth convergence path
				\item \textbf{Slow:} Each update requires O(N) computation
				\item \textbf{Memory intensive:} May need to store all data
			\end{itemize}
			\item \textbf{Advantages:}
			\begin{itemize}
				\item Guaranteed convergence for convex functions
				\item No random noise in gradient estimates
				\item Easier to analyze theoretically
			\end{itemize}
			\item \textbf{Disadvantages:}
			\begin{itemize}
				\item Impractical for large datasets ($N > 10^5$)
				\item Redundant computations
				\item Cannot escape local minima easily
			\end{itemize}
		\end{itemize}
		
		\vspace{0.05cm}
		\textcolor{myred}{\textbf{WHY OTHER OPTIONS ARE WRONG:}}
		\begin{itemize}
			\item \textcolor{myred}{A:} This is SGD (Stochastic Gradient Descent)
			\item \textcolor{myred}{C:} This is Mini-Batch Gradient Descent
			\item \textcolor{myred}{D:} This is random search
		\end{itemize}
		
		\textcolor{myred}{\textbf{EXAM TRAP:}} Batch GD uses \textcolor{myblue}{ALL data} for each update!
	\end{frame}
	
	\begin{frame}{Q22: Stochastic Gradient Descent - Tutorial}
		\fontsize{6.5}{7.5}\selectfont
		\textbf{What is Stochastic Gradient Descent (SGD)?}
		
		\vspace{0.05cm}
		\textcolor{myblue}{\textbf{ANSWER: B}} - Uses a single randomly selected data point per iteration.
		
		\vspace{0.05cm}
		\textbf{DETAILED EXPLANATION:}
		\begin{itemize}
			\item \textbf{Definition:} Computes gradient using one randomly chosen sample
			\begin{itemize}
				\item $\nabla L(\theta) \approx \nabla L_i(\theta)$ for random $i$
				\item One update per data point (or per mini-batch)
				\item Many updates per epoch
			\end{itemize}
			\item \textbf{Characteristics:}
			\begin{itemize}
				\item \textbf{Noisy gradient:} High variance in gradient estimates
				\item \textbf{Fast:} Each update is O(1) computation
				\item \textbf{Escapes local minima:} Noise helps explore the landscape
				\item \textbf{Converges to optimum:} With decreasing learning rate
			\end{itemize}
			\item \textbf{Advantages:}
			\begin{itemize}
				\item Scalable to huge datasets
				\item Can escape shallow local minima
				\item Online learning possible (streaming data)
			\end{itemize}
			\item \textbf{Disadvantages:}
			\begin{itemize}
				\item Noisy convergence (doesn't settle at exact minimum)
				\item Requires careful learning rate scheduling
				\item Harder to analyze theoretically
			\end{itemize}
		\end{itemize}
		
		\vspace{0.05cm}
		\textcolor{myred}{\textbf{WHY OTHER OPTIONS ARE WRONG:}}
		\begin{itemize}
			\item \textcolor{myred}{A:} This is Batch GD (uses all data)
			\item \textcolor{myred}{C:} This is Mini-Batch GD (uses subset)
			\item \textcolor{myred}{D:} This describes learning rate scheduling
		\end{itemize}
		
		\textcolor{myred}{\textbf{EXAM TRAP:}} SGD uses \textcolor{myblue}{ONE data point} per update!
	\end{frame}
	
	\begin{frame}{Q23: Mini-Batch Gradient Descent - Tutorial}
		\fontsize{6.5}{7.5}\selectfont
		\textbf{What is Mini-Batch Gradient Descent?}
		
		\vspace{0.05cm}
		\textcolor{myblue}{\textbf{ANSWER: C}} - Uses a small random subset (batch) of data per update.
		
		\vspace{0.05cm}
		\textbf{DETAILED EXPLANATION:}
		\begin{itemize}
			\item \textbf{Definition:} Computes gradient using a small batch of size $b$
			\begin{itemize}
				\item $\nabla L(\theta) \approx \frac{1}{b}\sum_{i \in \text{batch}} \nabla L_i(\theta)$
				\item Typical batch sizes: 32, 64, 128, 256
				\item One update per batch, multiple batches per epoch
			\end{itemize}
			\item \textbf{Characteristics:}
			\begin{itemize}
				\item \textbf{Balanced:} Compromise between Batch GD and SGD
				\item \textbf{Less noisy than SGD:} Averaging reduces variance
				\item \textbf{Faster than Batch GD:} Smaller computation per update
				\item \textbf{Better than SGD:} More stable convergence
			\end{itemize}
			\item \textbf{Advantages:}
			\begin{itemize}
				\item Efficient for GPU processing (batch matrix operations)
				\item Better gradient estimates than SGD
				\item Faster convergence than Batch GD
				\item Most common in deep learning practice
			\end{itemize}
			\item \textbf{Batch size choice:}
			\begin{itemize}
				\item Power of 2 for hardware efficiency (32, 64, 128)
				\item Too small: noisy like SGD
				\item Too large: slow like Batch GD
			\end{itemize}
		\end{itemize}
		
		\vspace{0.05cm}
		\textcolor{myred}{\textbf{WHY OTHER OPTIONS ARE WRONG:}}
		\begin{itemize}
			\item \textcolor{myred}{A:} This is Batch GD (all data)
			\item \textcolor{myred}{B:} This is SGD (one data point)
			\item \textcolor{myred}{D:} This describes a stopping criterion
		\end{itemize}
		
		\textcolor{myred}{\textbf{EXAM TRAP:}} Mini-Batch GD uses a \textcolor{myblue}{small random subset} of data!
	\end{frame}
	
	\begin{frame}{Q24: SGD Advantage - Tutorial}
		\fontsize{6.5}{7.5}\selectfont
		\textbf{Why is SGD often preferred over Batch GD for large datasets?}
		
		\vspace{0.05cm}
		\textcolor{myblue}{\textbf{ANSWER: B}} - Computationally cheaper per iteration (fewer data points).
		
		\vspace{0.05cm}
		\textbf{DETAILED EXPLANATION:}
		\begin{itemize}
			\item \textbf{Computational cost comparison:}
			\begin{itemize}
				\item Batch GD: $O(N)$ per update (entire dataset)
				\item SGD: $O(1)$ per update (one data point)
				\item Mini-batch: $O(b)$ per update (batch size $b$)
			\end{itemize}
			\item \textbf{Scaling to big data:}
			\begin{itemize}
				\item $N = 10^6$: Batch GD update requires 1M gradient computations
				\item SGD update requires 1 gradient computation
				\item With SGD, can make 1M updates in the time of one Batch update
			\end{itemize}
			\item \textbf{SGD convergence properties:}
			\begin{itemize}
				\item May require more iterations to converge
				\item But each iteration is much cheaper
				\item Total time often much less than Batch GD
			\end{itemize}
			\item \textbf{Additional benefits:}
			\begin{itemize}
				\item Can escape local minima due to noise
				\item Enables online learning
				\item Works with streaming data
			\end{itemize}
		\end{itemize}
		
		\vspace{0.05cm}
		\textcolor{myred}{\textbf{WHY OTHER OPTIONS ARE WRONG:}}
		\begin{itemize}
			\item \textcolor{myred}{A:} Not always - SGD converges noisily
			\item \textcolor{myred}{C:} False - SGD doesn't guarantee global minimum
			\item \textcolor{myred}{D:} False - SGD requires a loss function
		\end{itemize}
		
		\textcolor{myred}{\textbf{EXAM TRAP:}} SGD preferred for large datasets due to \textcolor{myblue}{lower computational cost per iteration}!
	\end{frame}
	
	\begin{frame}{Q25: Local Minima - Tutorial}
		\fontsize{6.5}{7.5}\selectfont
		\textbf{What is a local minimum in gradient descent?}
		
		\vspace{0.05cm}
		\textcolor{myblue}{\textbf{ANSWER: B}} - Gradient is zero but function value higher than at other nearby points.
		
		\vspace{0.05cm}
		\textbf{DETAILED EXPLANATION:}
		\begin{itemize}
			\item \textbf{Definition:}
			\begin{itemize}
				\item Point $\theta^*$ where $\nabla L(\theta^*) = 0$
				\item $L(\theta^*) < L(\theta)$ for all $\theta$ in neighborhood
				\item But not necessarily the lowest point overall
			\end{itemize}
			\item \textbf{Mathematical characterization:}
			\begin{itemize}
				\item First-order condition: $\nabla L(\theta^*) = 0$
				\item Second-order condition: Hessian $H(\theta^*)$ is positive semidefinite
				\item All eigenvalues of $H$ are $\geq 0$
			\end{itemize}
			\item \textbf{Why problematic:}
			\begin{itemize}
				\item Gradient descent stops (or slows) at local minima
				\item Cannot distinguish from global minimum
				\item May not find the best solution
			\end{itemize}
			\item \textbf{Local vs Global minimum:}
			\begin{itemize}
				\item Local minimum: best in local neighborhood
				\item Global minimum: best overall
				\item Convex functions: every local minimum is global
				\item Non-convex functions: multiple local minima possible
			\end{itemize}
		\end{itemize}
		
		\vspace{0.05cm}
		\textcolor{myred}{\textbf{WHY OTHER OPTIONS ARE WRONG:}}
		\begin{itemize}
			\item \textcolor{myred}{A:} This is global minimum (lowest overall)
			\item \textcolor{myred}{C:} Gradient is zero at minima, not infinite
			\item \textcolor{myred}{D:} This is a maximum (gradient zero, higher than neighbors)
		\end{itemize}
		
		\textcolor{myred}{\textbf{EXAM TRAP:}} Local minima are \textcolor{myred}{traps} where gradient descent can get stuck!
	\end{frame}
	
	\begin{frame}{Q26: Saddle Points - Tutorial}
		\fontsize{6.5}{7.5}\selectfont
		\textbf{What is a saddle point in gradient descent?}
		
		\vspace{0.05cm}
		\textcolor{myblue}{\textbf{ANSWER: B}} - Gradient zero but curves upward in some directions and downward in others.
		
		\vspace{0.05cm}
		\textbf{DETAILED EXPLANATION:}
		\begin{itemize}
			\item \textbf{Definition:}
			\begin{itemize}
				\item Point $\theta^*$ where $\nabla L(\theta^*) = 0$
				\item Not a local minimum or maximum
				\item Hessian has both positive and negative eigenvalues
			\end{itemize}
			\item \textbf{Mathematical characterization:}
			\begin{itemize}
				\item First-order: $\nabla L(\theta^*) = 0$
				\item Second-order: Hessian $H(\theta^*)$ is indefinite
				\item Some eigenvalues $> 0$ (upward curves)
				\item Some eigenvalues $< 0$ (downward curves)
			\end{itemize}
			\item \textbf{Why important:}
			\begin{itemize}
				\item More common than local minima in high dimensions
				\item Gradient descent slows down near saddle points
				\item Requires momentum or noise to escape
			\end{itemize}
			\item \textbf{Comparison to local minima:}
			\begin{itemize}
				\item Local minima: all eigenvalues positive
				\item Saddle points: mixed eigenvalues
				\item In high dimensions, saddle points dominate
			\end{itemize}
		\end{itemize}
		
		\vspace{0.05cm}
		\textcolor{myred}{\textbf{WHY OTHER OPTIONS ARE WRONG:}}
		\begin{itemize}
			\item \textcolor{myred}{A:} This describes a local minimum
			\item \textcolor{myred}{C:} Gradient is zero at saddle points
			\item \textcolor{myred}{D:} This is global minimum
		\end{itemize}
		
		\textcolor{myred}{\textbf{EXAM TRAP:}} Saddle points have \textcolor{myblue}{zero gradient} but are neither minima nor maxima!
	\end{frame}
	
	\begin{frame}{Q27: Momentum - Tutorial}
		\fontsize{6.5}{7.5}\selectfont
		\textbf{What is the purpose of momentum in gradient descent?}
		
		\vspace{0.05cm}
		\textcolor{myblue}{\textbf{ANSWER: B}} - Accelerate convergence and reduce oscillations by accumulating past gradients.
		
		\vspace{0.05cm}
		\textbf{DETAILED EXPLANATION:}
		\begin{itemize}
			\item \textbf{Definition:} Momentum adds fraction of previous update to current update
			\begin{itemize}
				\item $v_t = \beta v_{t-1} + (1-\beta)\nabla L(\theta_t)$
				\item $\theta_{t+1} = \theta_t - \alpha v_t$
				\item $\beta$ is momentum parameter (typically 0.9)
			\end{itemize}
			\item \textbf{How it works:}
			\begin{itemize}
				\item Accumulates gradient history
				\item Accelerates in consistent gradient directions
				\item Dampens oscillations in perpendicular directions
			\end{itemize}
			\item \textbf{Benefits:}
			\begin{itemize}
				\item \textbf{Faster convergence:} Smooths out zigzagging
				\item \textbf{Better local minima:} Can roll past shallow minima
				\item \textbf{Escape saddle points:} Built-up momentum pushes through
			\end{itemize}
			\item \textbf{Analogy:} Ball rolling down a hill
			\begin{itemize}
				\item Gains speed as it rolls
				\item Overcomes small bumps
				\item Doesn't stop at first flat spot
			\end{itemize}
		\end{itemize}
		
		\vspace{0.05cm}
		\textcolor{myred}{\textbf{WHY OTHER OPTIONS ARE WRONG:}}
		\begin{itemize}
			\item \textcolor{myred}{A:} Learning rate decay is separate concept
			\item \textcolor{myred}{C:} Batch size is separate concept
			\item \textcolor{myred}{D:} Initialization is separate step
		\end{itemize}
		
		\textcolor{myred}{\textbf{EXAM TRAP:}} Momentum uses \textcolor{myblue}{past gradients} to smooth updates!
	\end{frame}
	
	\begin{frame}{Q28: Adam Optimizer - Tutorial}
		\fontsize{6.5}{7.5}\selectfont
		\textbf{What is the Adam optimizer?}
		
		\vspace{0.05cm}
		\textcolor{myblue}{\textbf{ANSWER: B}} - Adaptive learning rate optimizer combining momentum and RMSprop.
		
		\vspace{0.05cm}
		\textbf{DETAILED EXPLANATION:}
		\begin{itemize}
			\item \textbf{Definition:} Adaptive Moment Estimation
			\begin{itemize}
				\item Combines advantages of two methods
				\item Momentum: tracks first moment (mean of gradients)
				\item RMSprop: tracks second moment (variance of gradients)
			\end{itemize}
			\item \textbf{Update formulas:}
			\begin{itemize}
				\item $m_t = \beta_1 m_{t-1} + (1-\beta_1)\nabla L(\theta_t)$
				\item $v_t = \beta_2 v_{t-1} + (1-\beta_2)(\nabla L(\theta_t))^2$
				\item $\hat{m}_t = m_t / (1-\beta_1^t)$ (bias correction)
				\item $\hat{v}_t = v_t / (1-\beta_2^t)$ (bias correction)
				\item $\theta_{t+1} = \theta_t - \alpha \hat{m}_t / (\sqrt{\hat{v}_t} + \varepsilon)$
			\end{itemize}
			\item \textbf{Key features:}
			\begin{itemize}
				\item Adaptive learning rate per parameter
				\item Handles sparse gradients well
				\item Requires minimal tuning
				\item Most popular optimizer in deep learning
			\end{itemize}
			\item \textbf{Default hyperparameters:}
			\begin{itemize}
				\item $\alpha = 0.001$ (learning rate)
				\item $\beta_1 = 0.9$ (momentum decay)
				\item $\beta_2 = 0.999$ (RMSprop decay)
				\item $\varepsilon = 10^{-8}$ (small constant for stability)
			\end{itemize}
		\end{itemize}
		
		\vspace{0.05cm}
		\textcolor{myred}{\textbf{WHY OTHER OPTIONS ARE WRONG:}}
		\begin{itemize}
			\item \textcolor{myred}{A:} Adam adapts learning rate dynamically
			\item \textcolor{myred}{C:} Adam uses first AND second moments
			\item \textcolor{myred}{D:} Adam is an optimizer, not initialization
		\end{itemize}
		
		\textcolor{myred}{\textbf{EXAM TRAP:}} Adam is \textcolor{myblue}{adaptive} - adjusts learning rates per parameter!
	\end{frame}
	
	\begin{frame}{Q29: Learning Rate Scheduling - Tutorial}
		\fontsize{6.5}{7.5}\selectfont
		\textbf{What is learning rate scheduling?}
		
		\vspace{0.05cm}
		\textcolor{myblue}{\textbf{ANSWER: B}} - Adjusts learning rate during training according to predefined schedule.
		
		\vspace{0.05cm}
		\textbf{DETAILED EXPLANATION:}
		\begin{itemize}
			\item \textbf{Definition:} Systematic reduction of learning rate over time
			\begin{itemize}
				\item Larger learning rates initially (exploration)
				\item Smaller learning rates later (fine-tuning)
			\end{itemize}
			\item \textbf{Why schedule is needed:}
			\begin{itemize}
				\item Large LR: make progress quickly, explore the space
				\item Small LR: settle into minimum, avoid oscillation
				\item Both are needed but at different times
			\end{itemize}
			\item \textbf{Common schedules:}
			\begin{enumerate}
				\item \textbf{Step decay:} Reduce by factor at fixed intervals
				\item \textbf{Exponential decay:} $\alpha_t = \alpha_0 \times \gamma^t$
				\item \textbf{Cosine annealing:} Smooth cyclical variation
				\item \textbf{Warmup:} Increase then decrease
			\end{enumerate}
			\item \textbf{Practical advice:}
			\begin{itemize}
				\item Start with larger learning rate
				\item Reduce when validation loss plateaus
				\item Different schedules work for different problems
			\end{itemize}
		\end{itemize}
		
		\vspace{0.05cm}
		\textcolor{myred}{\textbf{WHY OTHER OPTIONS ARE WRONG:}}
		\begin{itemize}
			\item \textcolor{myred}{A:} Random changes are not a "schedule"
			\item \textcolor{myred}{C:} Constant LR is the opposite of scheduling
			\item \textcolor{myred}{D:} Batch size is a different hyperparameter
		\end{itemize}
		
		\textcolor{myred}{\textbf{EXAM TRAP:}} Learning rate schedules \textcolor{myblue}{adjust learning rate over time}!
	\end{frame}
	
	\begin{frame}{Q30: Step Decay - Tutorial}
		\fontsize{6.5}{7.5}\selectfont
		\textbf{What is step decay in learning rate scheduling?}
		
		\vspace{0.05cm}
		\textcolor{myblue}{\textbf{ANSWER: B}} - Learning rate decreases by factor at specific intervals.
		
		\vspace{0.05cm}
		\textbf{DETAILED EXPLANATION:}
		\begin{itemize}
			\item \textbf{Definition:} Piecewise constant reduction
			\begin{itemize}
				\item $\alpha_t = \alpha_0 \times \gamma^{\lfloor t/T \rfloor}$
				\item $\gamma$: decay factor (e.g., 0.1 or 0.5)
				\item $T$: interval size (e.g., every 10 epochs)
			\end{itemize}
			\item \textbf{Example:}
			\begin{itemize}
				\item $\alpha_0 = 0.01$, $\gamma = 0.1$, $T = 10$
				\item Epochs 0-9: $\alpha = 0.01$
				\item Epochs 10-19: $\alpha = 0.001$
				\item Epochs 20-29: $\alpha = 0.0001$
			\end{itemize}
			\item \textbf{Advantages:}
			\begin{itemize}
				\item Simple and easy to implement
				\item Predictable behavior
				\item Works well for many problems
			\end{itemize}
			\item \textbf{Disadvantages:}
			\begin{itemize}
				\item Abrupt changes in learning rate
				\item Requires choosing interval and factor
				\item Not adaptive to data
			\end{itemize}
			\item \textbf{When to use:}
			\begin{itemize}
				\item When training loss plateaus
				\item When model is stable
				\item Common in practice
			\end{itemize}
		\end{itemize}
		
		\vspace{0.05cm}
		\textcolor{myred}{\textbf{WHY OTHER OPTIONS ARE WRONG:}}
		\begin{itemize}
			\item \textcolor{myred}{A:} Continuous decrease is exponential decay
			\item \textcolor{myred}{C:} Increasing LR is a different strategy
			\item \textcolor{myred}{D:} Random changes are not step decay
		\end{itemize}
		
		\textcolor{myred}{\textbf{EXAM TRAP:}} Step decay reduces learning rate \textcolor{myblue}{at specific intervals}!
	\end{frame}
	
	\begin{frame}{Q31: Exponential Decay - Tutorial}
		\fontsize{6.5}{7.5}\selectfont
		\textbf{What is exponential decay in learning rate scheduling?}
		
		\vspace{0.05cm}
		\textcolor{myblue}{\textbf{ANSWER: A}} - $\alpha_t = \alpha_0 \times \gamma^{t}$
		
		\vspace{0.05cm}
		\textbf{DETAILED EXPLANATION:}
		\begin{itemize}
			\item \textbf{Definition:} Continuous multiplicative decay
			\begin{itemize}
				\item $\alpha_t = \alpha_0 \times \gamma^{t}$
				\item $\gamma \in (0, 1)$: decay rate
				\item $\alpha_0$: initial learning rate
			\end{itemize}
			\item \textbf{Example:}
			\begin{itemize}
				\item $\alpha_0 = 0.01$, $\gamma = 0.95$
				\item Epoch 0: $\alpha = 0.01$
				\item Epoch 10: $\alpha = 0.01 \times 0.95^{10} \approx 0.006$
				\item Epoch 100: $\alpha = 0.01 \times 0.95^{100} \approx 0.00006$
			\end{itemize}
			\item \textbf{Properties:}
			\begin{itemize}
				\item \textbf{Smooth decay:} No abrupt changes
				\item \textbf{Consistent:} Same relative decrease each step
				\item \textbf{Long tail:} Never reaches zero (asymptotically)
			\end{itemize}
			\item \textbf{Choosing $\gamma$:}
			\begin{itemize}
				\item $\gamma = 0.9$: faster decay
				\item $\gamma = 0.99$: slower decay
				\item $\gamma = 0.95$: common default
			\end{itemize}
			\item \textbf{Comparison to step decay:}
			\begin{itemize}
				\item Step decay: discrete jumps
				\item Exponential decay: smooth continuous decrease
			\end{itemize}
		\end{itemize}
		
		\vspace{0.05cm}
		\textcolor{myred}{\textbf{WHY OTHER OPTIONS ARE WRONG:}}
		\begin{itemize}
			\item \textcolor{myred}{B:} This is inverse time decay ($\alpha_t = \alpha_0 / t$)
			\item \textcolor{myred}{C:} This is linear decay
			\item \textcolor{myred}{D:} This is constant learning rate
		\end{itemize}
		
		\textcolor{myred}{\textbf{EXAM TRAP:}} Exponential decay: $\alpha_t = \alpha_0 \times \gamma^{t}$!
	\end{frame}
	
	\begin{frame}{Q32: Convergence Definition - Tutorial}
		\fontsize{6.5}{7.5}\selectfont
		\textbf{What does convergence mean in gradient descent?}
		
		\vspace{0.05cm}
		\textcolor{myblue}{\textbf{ANSWER: C}} - Loss has stopped decreasing significantly (or gradient near zero).
		
		\vspace{0.05cm}
		\textbf{DETAILED EXPLANATION:}
		\begin{itemize}
			\item \textbf{Definition:} Algorithm has reached a stationary point
			\begin{itemize}
				\item Gradient norm $\|\nabla L(\theta)\| < \varepsilon$
				\item Loss change $\Delta L < \delta$
				\item Parameter change $\|\theta_t - \theta_{t-1}\| < \tau$
			\end{itemize}
			\item \textbf{Convergence criteria:}
			\begin{enumerate}
				\item \textbf{Gradient norm:} $\|\nabla L(\theta_t)\| < 10^{-5}$
				\item \textbf{Loss change:} $|L(\theta_t) - L(\theta_{t-1})| < 10^{-6}$
				\item \textbf{Parameter change:} $\|\theta_t - \theta_{t-1}\| < 10^{-5}$
				\item \textbf{Validation loss:} Stopping when validation loss increases
			\end{enumerate}
			\item \textbf{Convergence vs global minimum:}
			\begin{itemize}
				\item Convergence $\neq$ global minimum
				\item Could be local minimum or saddle point
				\item Function is convex: convergence $\Rightarrow$ global minimum
			\end{itemize}
			\item \textbf{Stopping vs convergence:}
			\begin{itemize}
				\item Max iterations: stopping criterion, not convergence
				\item Convergence: based on progress, not iteration count
			\end{itemize}
		\end{itemize}
		
		\vspace{0.05cm}
		\textcolor{myred}{\textbf{WHY OTHER OPTIONS ARE WRONG:}}
		\begin{itemize}
			\item \textcolor{myred}{A:} Ideal case but not guaranteed
			\item \textcolor{myred}{B:} Could be due to zero learning rate
			\item \textcolor{myred}{D:} This is a stopping criterion, not convergence
		\end{itemize}
		
		\textcolor{myred}{\textbf{EXAM TRAP:}} Convergence means \textcolor{myblue}{loss has stopped decreasing} significantly!
	\end{frame}
	
	\begin{frame}{Q33: Loss Landscape - Tutorial}
		\fontsize{6.5}{7.5}\selectfont
		\textbf{In high-dimensional spaces, which is more common: local minima or saddle points?}
		
		\vspace{0.05cm}
		\textcolor{myblue}{\textbf{ANSWER: B}} - Saddle points are more common.
		
		\vspace{0.05cm}
		\textbf{DETAILED EXPLANATION:}
		\begin{itemize}
			\item \textbf{Theoretical result:} Saddle points dominate in high dimensions
			\begin{itemize}
				\item For a random function in $d$ dimensions
				\item Probability of a critical point being local minimum: $2^{-d}$
				\item Probability of being saddle point: $1 - 2^{-d}$
				\item For $d = 100$: $P(\text{local min}) \approx 10^{-30}$
			\end{itemize}
			\item \textbf{Intuition:}
			\begin{itemize}
				\item Local minimum requires all directions to be upward
				\item Saddle point has some up and some down directions
				\item More likely to have mixed directions in high dimensions
			\end{itemize}
			\item \textbf{Implications for deep learning:}
			\begin{itemize}
				\item Local minima are not the main problem
				\item Saddle points are more common
				\item Gradient descent can escape saddle points
				\item Momentum and noise help escape
			\end{itemize}
			\item \textbf{Recent research findings:}
			\begin{itemize}
				\item Deep networks have many saddle points
				\item Most local minima are nearly global
				\item The loss landscape is surprisingly well-behaved
			\end{itemize}
		\end{itemize}
		
		\vspace{0.05cm}
		\textcolor{myred}{\textbf{WHY OTHER OPTIONS ARE WRONG:}}
		\begin{itemize}
			\item \textcolor{myred}{A:} False - local minima are rare in high dimensions
			\item \textcolor{myred}{C:} False - no guarantee of global minimum
			\item \textcolor{myred}{D:} False - both exist but saddle points more common
		\end{itemize}
		
		\textcolor{myred}{\textbf{EXAM TRAP:}} Saddle points are \textcolor{myblue}{more common} than local minima in high dimensions!
	\end{frame}
	
	\begin{frame}{Q34: Vanishing Gradients - Tutorial}
		\fontsize{6.5}{7.5}\selectfont
		\textbf{What is the vanishing gradient problem?}
		
		\vspace{0.05cm}
		\textcolor{myblue}{\textbf{ANSWER: B}} - Gradients become very small during backpropagation, preventing early layers from learning.
		
		\vspace{0.05cm}
		\textbf{DETAILED EXPLANATION:}
		\begin{itemize}
			\item \textbf{Definition:} Gradients shrink exponentially as they propagate backward
			\begin{itemize}
				\item Early layers receive very small gradients
				\item Weights in early layers barely update
				\item Network stops learning effectively
			\end{itemize}
			\item \textbf{Causes:}
			\begin{itemize}
				\item Activation functions with small derivatives
				\item \textbf{Sigmoid:} derivative $\sigma(x)(1-\sigma(x)) \leq 0.25$
				\item \textbf{tanh:} derivative $\text{sech}^2(x) \to 0$ for large $|x|$
				\item Deep networks (many layers multiply small values)
			\end{itemize}
			\item \textbf{Signs of vanishing gradients:}
			\begin{itemize}
				\item Early layers have very small gradient norms
				\item Training loss doesn't decrease
				\item Early layers look like random initialization
			\end{itemize}
			\item \textbf{Solutions:}
			\begin{itemize}
				\item ReLU activation (derivative 1 for positive inputs)
				\item Residual connections (skip connections)
				\item Batch normalization
				\item Careful weight initialization (Xavier, He)
			\end{itemize}
		\end{itemize}
		
		\vspace{0.05cm}
		\textcolor{myred}{\textbf{WHY OTHER OPTIONS ARE WRONG:}}
		\begin{itemize}
			\item \textcolor{myred}{A:} This describes exploding gradients
			\item \textcolor{myred}{C:} This describes oscillations
			\item \textcolor{myred}{D:} Extreme case, not the general problem
		\end{itemize}
		
		\textcolor{myred}{\textbf{EXAM TRAP:}} Vanishing gradients = \textcolor{myblue}{very small gradients} in early layers!
	\end{frame}
	
	\begin{frame}{Q35: Exploding Gradients - Tutorial}
		\fontsize{6.5}{7.5}\selectfont
		\textbf{What is the exploding gradient problem?}
		
		\vspace{0.05cm}
		\textcolor{myblue}{\textbf{ANSWER: B}} - Gradients become very large during backpropagation, causing unstable updates.
		
		\vspace{0.05cm}
		\textbf{DETAILED EXPLANATION:}
		\begin{itemize}
			\item \textbf{Definition:} Gradients grow exponentially as they propagate backward
			\begin{itemize}
				\item Large gradients lead to large parameter updates
				\item Parameters become very large
				\item Training becomes unstable or diverges
			\end{itemize}
			\item \textbf{Causes:}
			\begin{itemize}
				\item Large initial weights
				\item Poor weight initialization
				\item Long sequences (RNNs)
				\item Repeated multiplication of values $> 1$
			\end{itemize}
			\item \textbf{Signs of exploding gradients:}
			\begin{itemize}
				\item Gradients become extremely large ($> 10^6$)
				\item Loss becomes NaN (not a number)
				\item Parameters become extremely large
				\item Training diverges (loss increases)
			\end{itemize}
			\item \textbf{Solutions:}
			\begin{itemize}
				\item \textbf{Gradient clipping:} Cap gradients at threshold
				\begin{itemize}
					\item Norm clipping: $\nabla L \leftarrow \frac{\nabla L}{\|\nabla L\|} \times \text{threshold}$
					\item Value clipping: $\nabla L \leftarrow \min(\max(\nabla L, -\text{threshold}), \text{threshold})$
				\end{itemize}
				\item Careful weight initialization
				\item Smaller learning rates
				\item L2 regularization (weight decay)
			\end{itemize}
		\end{itemize}
		
		\vspace{0.05cm}
		\textcolor{myred}{\textbf{WHY OTHER OPTIONS ARE WRONG:}}
		\begin{itemize}
			\item \textcolor{myred}{A:} This describes vanishing gradients
			\item \textcolor{myred}{C:} This is opposite (zero gradients)
			\item \textcolor{myred}{D:} This describes a different phenomenon
		\end{itemize}
		
		\textcolor{myred}{\textbf{EXAM TRAP:}} Exploding gradients = \textcolor{myblue}{very large gradients} causing instability!
	\end{frame}
	
	% ============================================================
	% SECTION 3: BACKPROPAGATION - QUESTIONS 36-45
	% ============================================================
	
	\begin{frame}{Q36: Backpropagation Definition - Tutorial}
		\fontsize{6.5}{7.5}\selectfont
		\textbf{What is backpropagation in neural networks?}
		
		\vspace{0.05cm}
		\textcolor{myblue}{\textbf{ANSWER: B}} - Algorithm that computes gradients of loss with respect to all weights using chain rule.
		
		\vspace{0.05cm}
		\textbf{DETAILED EXPLANATION:}
		\begin{itemize}
			\item \textbf{Definition:} Efficient algorithm for computing gradients
			\begin{itemize}
				\item Uses chain rule of calculus
				\item Propagates errors backward from output to input
				\item Computes $\frac{\partial L}{\partial w_i}$ for every weight
			\end{itemize}
			\item \textbf{Algorithm steps:}
			\begin{enumerate}
				\item \textbf{Forward pass:} Compute predictions and loss
				\item \textbf{Backward pass:} Compute gradients of loss
				\item \textbf{Update:} Use gradients to update weights
			\end{enumerate}
			\item \textbf{Key insight:}
			\begin{itemize}
				\item Reuses computations efficiently
				\item Each gradient computed once
				\item O(N) complexity for N parameters
			\end{itemize}
			\item \textbf{Historical context:}
			\begin{itemize}
				\item Invented independently multiple times
				\item Popularized by Rumelhart, Hinton, Williams (1986)
				\item Made deep learning possible
			\end{itemize}
		\end{itemize}
		
		\vspace{0.05cm}
		\textcolor{myred}{\textbf{WHY OTHER OPTIONS ARE WRONG:}}
		\begin{itemize}
			\item \textcolor{myred}{A:} Initialization is a separate step
			\item \textcolor{myred}{C:} Architecture selection is separate
			\item \textcolor{myred}{D:} Activation function selection is separate
		\end{itemize}
		
		\textcolor{myred}{\textbf{EXAM TRAP:}} Backpropagation uses \textcolor{myblue}{chain rule} to compute gradients!
	\end{frame}
	
	\begin{frame}{Q37: Chain Rule - Tutorial}
		\fontsize{6.5}{7.5}\selectfont
		\textbf{Which mathematical rule is the foundation of backpropagation?}
		
		\vspace{0.05cm}
		\textcolor{myblue}{\textbf{ANSWER: C}} - The chain rule.
		
		\vspace{0.05cm}
		\textbf{DETAILED EXPLANATION:}
		\begin{itemize}
			\item \textbf{Chain rule in calculus:}
			\begin{itemize}
				\item For composite functions: $\frac{d}{dx} f(g(x)) = f'(g(x)) g'(x)$
				\item Neural networks are compositions of many functions
				\item Loss $L = f_N(f_{N-1}(... f_1(x)...))$
			\end{itemize}
			\item \textbf{Backpropagation application:}
			\begin{itemize}
				\item $\frac{\partial L}{\partial w_{ij}^{(l)}} = \frac{\partial L}{\partial z_i^{(l+1)}} \cdot \frac{\partial z_i^{(l+1)}}{\partial z_j^{(l)}} \cdot ...$
				\item Each layer contributes one factor
				\item Gradients are multiplied backward
			\end{itemize}
			\item \textbf{Example:}
			\begin{itemize}
				\item $z = w_2 \cdot \sigma(w_1 x)$
				\item $\frac{\partial L}{\partial w_1} = \frac{\partial L}{\partial z} \cdot \frac{\partial z}{\partial a} \cdot \frac{\partial a}{\partial w_1}$
				\item $a = \sigma(w_1 x)$
			\end{itemize}
			\item \textbf{Why not other rules:}
			\begin{itemize}
				\item Product rule: for products, not compositions
				\item Quotient rule: for ratios
				\item Power rule: for powers
			\end{itemize}
		\end{itemize}
		
		\vspace{0.05cm}
		\textcolor{myred}{\textbf{WHY OTHER OPTIONS ARE WRONG:}}
		\begin{itemize}
			\item \textcolor{myred}{A:} Product rule applies to products of functions
			\item \textcolor{myred}{B:} Quotient rule applies to ratios
			\item \textcolor{myred}{D:} Power rule applies to powers
		\end{itemize}
		
		\textcolor{myred}{\textbf{EXAM TRAP:}} Backpropagation = \textcolor{myblue}{chain rule} applied to neural networks!
	\end{frame}
	
	\begin{frame}{Q38: Forward Pass - Tutorial}
		\fontsize{6.5}{7.5}\selectfont
		\textbf{What happens during forward pass in backpropagation?}
		
		\vspace{0.05cm}
		\textcolor{myblue}{\textbf{ANSWER: B}} - Input data passed through network to compute output and loss.
		
		\vspace{0.05cm}
		\textbf{DETAILED EXPLANATION:}
		\begin{itemize}
			\item \textbf{Forward pass steps:}
			\begin{enumerate}
				\item Input data enters the network
				\item Passed through hidden layers
				\item Each layer applies: $z^{(l)} = W^{(l)} a^{(l-1)} + b^{(l)}$
				\item Activation: $a^{(l)} = \sigma(z^{(l)})$
				\item Output layer produces predictions $\hat{y}$
				\item Loss $L$ is computed: $L = \text{Loss}(y, \hat{y})$
			\end{enumerate}
			\item \textbf{What is stored:}
			\begin{itemize}
				\item Inputs to each layer $z^{(l)}$
				\item Activations $a^{(l)}$
				\item Loss value
				\item Used in backward pass for gradient computation
			\end{itemize}
			\item \textbf{Complexity:}
			\begin{itemize}
				\item One forward pass per batch
				\item O(N) operations where N = number of parameters
				\item Must be computed before backward pass
			\end{itemize}
		\end{itemize}
		
		\vspace{0.05cm}
		\textcolor{myred}{\textbf{WHY OTHER OPTIONS ARE WRONG:}}
		\begin{itemize}
			\item \textcolor{myred}{A:} This is backward pass (gradients and updates)
			\item \textcolor{myred}{C:} Optimization happens over many iterations
			\item \textcolor{myred}{D:} Initialization happens before training
		\end{itemize}
		
		\textcolor{myred}{\textbf{EXAM TRAP:}} Forward pass computes \textcolor{myblue}{output and loss}; backward pass computes gradients!
	\end{frame}
	
	\begin{frame}{Q39: Backward Pass - Tutorial}
		\fontsize{6.5}{7.5}\selectfont
		\textbf{What happens during backward pass in backpropagation?}
		
		\vspace{0.05cm}
		\textcolor{myblue}{\textbf{ANSWER: B}} - Gradients of loss with respect to each weight computed using chain rule.
		
		\vspace{0.05cm}
		\textbf{DETAILED EXPLANATION:}
		\begin{itemize}
			\item \textbf{Backward pass steps:}
			\begin{enumerate}
				\item Start from output layer
				\item Compute $\delta^{(L)} = \frac{\partial L}{\partial a^{(L)}} \odot \sigma'(z^{(L)})$
				\item Propagate backward: $\delta^{(l)} = (W^{(l+1)})^T \delta^{(l+1)} \odot \sigma'(z^{(l)})$
				\item Compute weight gradients: $\frac{\partial L}{\partial W^{(l)}} = \delta^{(l)} (a^{(l-1)})^T$
				\item Compute bias gradients: $\frac{\partial L}{\partial b^{(l)}} = \delta^{(l)}$
			\end{enumerate}
			\item \textbf{What is used:}
			\begin{itemize}
				\item Activations from forward pass
				\item Loss gradient
				\item Chain rule applied layer by layer
			\end{itemize}
			\item \textbf{Complexity:}
			\begin{itemize}
				\item One backward pass per batch
				\item O(N) operations
				\item Computes gradients for all parameters
			\end{itemize}
		\end{itemize}
		
		\vspace{0.05cm}
		\textcolor{myred}{\textbf{WHY OTHER OPTIONS ARE WRONG:}}
		\begin{itemize}
			\item \textcolor{myred}{A:} This describes forward pass
			\item \textcolor{myred}{C:} Learning rate adjustment is separate
			\item \textcolor{myred}{D:} Loss evaluation happens during forward pass
		\end{itemize}
		
		\textcolor{myred}{\textbf{EXAM TRAP:}} Backward pass computes \textcolor{myblue}{gradients} of all weights!
	\end{frame}
	
	\begin{frame}{Q40: Weight Update - Tutorial}
		\fontsize{6.5}{7.5}\selectfont
		\textbf{After computing gradients, how are weights updated?}
		
		\vspace{0.05cm}
		\textcolor{myblue}{\textbf{ANSWER: B}} - $w_{new} = w_{old} - \alpha \nabla L(w)$
		
		\vspace{0.05cm}
		\textbf{DETAILED EXPLANATION:}
		\begin{itemize}
			\item \textbf{Update rule:}
			\begin{itemize}
				\item $w_{new} = w_{old} - \alpha \nabla L(w)$
				\item $\alpha$: learning rate
				\item $\nabla L(w)$: gradient computed by backpropagation
			\end{itemize}
			\item \textbf{Why this form:}
			\begin{itemize}
				\item Gradient points uphill (direction of increasing loss)
				\item Negative gradient points downhill (decreasing loss)
				\item Move in direction of decreasing loss
			\end{itemize}
			\item \textbf{Example for simple network:}
			\begin{itemize}
				\item $w_{new} = w_{old} - \alpha \frac{\partial L}{\partial w}$
				\item Applied to every weight and bias
				\item Update after each batch (or epoch)
			\end{itemize}
			\item \textbf{Variations:}
			\begin{itemize}
				\item Momentum: include previous updates
				\item Adam: adaptive learning rates
				\item Weight decay: add regularization
			\end{itemize}
		\end{itemize}
		
		\vspace{0.05cm}
		\textcolor{myred}{\textbf{WHY OTHER OPTIONS ARE WRONG:}}
		\begin{itemize}
			\item \textcolor{myred}{A:} This is gradient ascent (increases loss)
			\item \textcolor{myred}{C:} Ignores learning rate and old weights
			\item \textcolor{myred}{D:} Multiplication alone ignores gradient
		\end{itemize}
		
		\textcolor{myred}{\textbf{EXAM TRAP:}} Weights updated by \textcolor{myblue}{subtracting} gradient $\times$ learning rate!
	\end{frame}
	
	\begin{frame}{Q41: Computational Graph - Tutorial}
		\fontsize{6.5}{7.5}\selectfont
		\textbf{What is a computational graph in backpropagation?}
		
		\vspace{0.05cm}
		\textcolor{myblue}{\textbf{ANSWER: B}} - Graph representing forward computation as DAG enabling efficient gradient computation.
		
		\vspace{0.05cm}
		\textbf{DETAILED EXPLANATION:}
		\begin{itemize}
			\item \textbf{Definition:} Directed Acyclic Graph (DAG) of operations
			\begin{itemize}
				\item Nodes: operations (add, multiply, activation)
				\item Edges: data flow (tensors)
				\item Represents the entire computation
			\end{itemize}
			\item \textbf{Forward pass in graph:}
			\begin{itemize}
				\item Input nodes $x_1, x_2, ...$
				\item Operation nodes: $+$, $\times$, $\sigma$, ...
				\item Output: loss value
				\item Traverse forward from inputs to outputs
			\end{itemize}
			\item \textbf{Backward pass in graph:}
			\begin{itemize}
				\item Start from loss node
				\item Traverse backward using chain rule
				\item Each node computes gradient of its output
				\item Gradients propagate backward
			\end{itemize}
			\item \textbf{Frameworks using this:}
			\begin{itemize}
				\item TensorFlow (static graphs)
				\item PyTorch (dynamic graphs)
				\item JAX, Theano, etc.
			\end{itemize}
		\end{itemize}
		
		\vspace{0.05cm}
		\textcolor{myred}{\textbf{WHY OTHER OPTIONS ARE WRONG:}}
		\begin{itemize}
			\item \textcolor{myred}{A:} This is architecture diagram
			\item \textcolor{myred}{C:} This is data distribution
			\item \textcolor{myred}{D:} This is loss function plot
		\end{itemize}
		
		\textcolor{myred}{\textbf{EXAM TRAP:}} Computational graphs enable \textcolor{myblue}{automatic differentiation} and backpropagation!
	\end{frame}
	
	\begin{frame}{Q42: Automatic Differentiation - Tutorial}
		\fontsize{6.5}{7.5}\selectfont
		\textbf{What is automatic differentiation in deep learning frameworks?}
		
		\vspace{0.05cm}
		\textcolor{myblue}{\textbf{ANSWER: B}} - Automatic computation of gradients using chain rule on computational graphs.
		
		\vspace{0.05cm}
		\textbf{DETAILED EXPLANATION:}
		\begin{itemize}
			\item \textbf{Definition:} "Autodiff" = automatic gradient computation
			\begin{itemize}
				\item Uses computational graph
				\item Applies chain rule automatically
				\item Framework handles the details
			\end{itemize}
			\item \textbf{How it works:}
			\begin{enumerate}
				\item Build computational graph
				\item Forward pass computes outputs
				\item Backward pass computes gradients
				\item Returns gradients for all parameters
			\end{enumerate}
			\item \textbf{Modes of autodiff:}
			\begin{itemize}
				\item \textbf{Forward mode:} Compute derivatives with respect to inputs
				\item \textbf{Reverse mode:} Compute derivatives with respect to outputs (backpropagation)
				\item Reverse mode is more efficient for many parameters
			\end{itemize}
			\item \textbf{Advantages:}
			\begin{itemize}
				\item No manual derivative derivation
				\item Less error-prone
				\item Works for complex architectures
				\item Enables rapid prototyping
			\end{itemize}
		\end{itemize}
		
		\vspace{0.05cm}
		\textcolor{myred}{\textbf{WHY OTHER OPTIONS ARE WRONG:}}
		\begin{itemize}
			\item \textcolor{myred}{A:} Learning rate selection is separate
			\item \textcolor{myred}{C:} Architecture selection is separate
			\item \textcolor{myred}{D:} Data preprocessing is separate
		\end{itemize}
		
		\textcolor{myred}{\textbf{EXAM TRAP:}} Automatic differentiation is how frameworks implement \textcolor{myblue}{backpropagation}!
	\end{frame}
	
	\begin{frame}{Q43: Memory Requirements - Tutorial}
		\fontsize{6.5}{7.5}\selectfont
		\textbf{What is a key memory requirement for backpropagation?}
		
		\vspace{0.05cm}
		\textcolor{myblue}{\textbf{ANSWER: A}} - Requires storing all intermediate activations from forward pass.
		
		\vspace{0.05cm}
		\textbf{DETAILED EXPLANATION:}
		\begin{itemize}
			\item \textbf{Why activations must be stored:}
			\begin{itemize}
				\item Backward pass needs them to compute gradients
				\item $\frac{\partial L}{\partial w^{(l)}} = \delta^{(l)} (a^{(l-1)})^T$
				\item $a^{(l-1)}$ comes from forward pass
			\end{itemize}
			\item \textbf{Memory breakdown:}
			\begin{itemize}
				\item Model parameters: $O(p)$ memory
				\item Activations: $O(N \times \text{batch\_size} \times \text{layers})$
				\item Gradients: $O(p)$ memory
				\item Total: often dominated by activations
			\end{itemize}
			\item \textbf{Memory optimization techniques:}
			\begin{itemize}
				\item Checkpointing: recompute activations when needed
				\item Mixed precision: use lower precision storage
				\item Gradient checkpointing: trade compute for memory
				\item Smaller batch sizes
			\end{itemize}
			\item \textbf{Why this matters:}
			\begin{itemize}
				\item Limits batch size
				\item Limits model size
				\item Affects training speed
			\end{itemize}
		\end{itemize}
		
		\vspace{0.05cm}
		\textcolor{myred}{\textbf{WHY OTHER OPTIONS ARE WRONG:}}
		\begin{itemize}
			\item \textcolor{myred}{B:} False - activations are needed
			\item \textcolor{myred}{C:} False - final loss alone is insufficient
			\item \textcolor{myred}{D:} False - gradients are computed during backward pass
		\end{itemize}
		
		\textcolor{myred}{\textbf{EXAM TRAP:}} Backpropagation requires storing \textcolor{myblue}{all intermediate activations}!
	\end{frame}
	
	\begin{frame}{Q44: Training Loop Order - Tutorial}
		\fontsize{6.5}{7.5}\selectfont
		\textbf{What is correct order of operations in training loop?}
		
		\vspace{0.05cm}
		\textcolor{myblue}{\textbf{ANSWER: A}} - Forward pass $\to$ Compute loss $\to$ Backward pass $\to$ Update weights
		
		\vspace{0.05cm}
		\textbf{DETAILED EXPLANATION:}
		\begin{itemize}
			\item \textbf{Step-by-step breakdown:}
			\begin{enumerate}
				\item \textbf{Forward pass:}
				\begin{itemize}
					\item Input data through network
					\item Compute predictions $\hat{y}$
					\item Store activations
				\end{itemize}
				\item \textbf{Compute loss:}
				\begin{itemize}
					\item Compare $\hat{y}$ to actual $y$
					\item $L = \text{Loss}(y, \hat{y})$
				\end{itemize}
				\item \textbf{Backward pass:}
				\begin{itemize}
					\item Compute $\frac{\partial L}{\partial w}$ for all weights
					\item Use chain rule
					\item Propagate gradients backward
				\end{itemize}
				\item \textbf{Update weights:}
				\begin{itemize}
					\item $w_{new} = w_{old} - \alpha \nabla L(w)$
					\item Apply gradient descent step
				\end{itemize}
			\end{enumerate}
			\item \textbf{Why this order:}
			\begin{itemize}
				\item Need predictions before loss
				\item Need loss before gradients
				\item Need gradients before updates
			\end{itemize}
		\end{itemize}
		
		\vspace{0.05cm}
		\textcolor{myred}{\textbf{WHY OTHER OPTIONS ARE WRONG:}}
		\begin{itemize}
			\item \textcolor{myred}{B:} Backward pass cannot happen before forward pass
			\item \textcolor{myred}{C:} Loss cannot be computed before forward pass
			\item \textcolor{myred}{D:} Updates cannot happen before gradients
		\end{itemize}
		
		\textcolor{myred}{\textbf{EXAM TRAP:}} Forward $\to$ Loss $\to$ Backward $\to$ Update!
	\end{frame}
	
	\begin{frame}{Q45: Backpropagation Summary - Tutorial}
		\fontsize{6.5}{7.5}\selectfont
		\textbf{Which statement about backpropagation is FALSE?}
		
		\vspace{0.05cm}
		\textcolor{myblue}{\textbf{ANSWER: D}} - Backpropagation always finds the global minimum of the loss function.
		
		\vspace{0.05cm}
		\textbf{DETAILED EXPLANATION:}
		\begin{itemize}
			\item \textbf{Why D is false:}
			\begin{itemize}
				\item No guarantee of global minimum
				\item Can get stuck in local minima
				\item Can get trapped in saddle points
				\item Only guaranteed for convex functions
				\item Neural networks are non-convex
			\end{itemize}
			\item \textbf{What backpropagation actually does:}
			\begin{itemize}
				\item Computes gradients efficiently
				\item Enables gradient descent
				\item Finds a local minimum (often good enough)
				\item Converges to stationary point
			\end{itemize}
			\item \textbf{Why other statements are true:}
			\begin{itemize}
				\item \textcolor{mygreen}{A:} True - computes gradients for all weights
				\item \textcolor{mygreen}{B:} True - uses chain rule
				\item \textcolor{mygreen}{C:} True - requires storing activations
			\end{itemize}
		\end{itemize}
		
		\vspace{0.05cm}
		\textcolor{myred}{\textbf{WHY OTHER OPTIONS ARE WRONG:}}
		\begin{itemize}
			\item \textcolor{myred}{A:} This is exactly what backpropagation does
			\item \textcolor{myred}{B:} This is the mathematical foundation
			\item \textcolor{myred}{C:} This is the memory requirement
		\end{itemize}
		
		\textcolor{myred}{\textbf{EXAM TRAP:}} Backpropagation does \textcolor{myred}{NOT} guarantee finding the global minimum!
	\end{frame}
	
	% ============================================================
	% SECTION 4: BIAS-VARIANCE TRADEOFF - QUESTIONS 46-60
	% ============================================================
	
	\begin{frame}{Q46: Bias-Variance Tradeoff - Tutorial}
		\fontsize{6.5}{7.5}\selectfont
		\textbf{What is the bias-variance tradeoff in supervised learning?}
		
		\vspace{0.05cm}
		\textcolor{myblue}{\textbf{ANSWER: B}} - Tradeoff between model's ability to fit training data (bias) and generalize to new data (variance).
		
		\vspace{0.05cm}
		\textbf{DETAILED EXPLANATION:}
		\begin{itemize}
			\item \textbf{Mathematical decomposition:}
			\begin{itemize}
				\item $E[(Y - \hat{f}(X))^2] = \text{Bias}^2 + \text{Variance} + \text{Irreducible Error}$
				\item Bias$^2$: $[E[\hat{f}(X)] - f(X)]^2$
				\item Variance: $E[(\hat{f}(X) - E[\hat{f}(X)])^2]$
				\item Irreducible: $\sigma^2$ (noise in data)
			\end{itemize}
			\item \textbf{Intuition:}
			\begin{itemize}
				\item \textbf{Bias:} How wrong the model is on average
				\item \textbf{Variance:} How much the model changes with different training data
				\item Tradeoff: Can't reduce both simultaneously
			\end{itemize}
			\item \textbf{Model complexity relationship:}
			\begin{itemize}
				\item Simple models: High bias, Low variance
				\item Complex models: Low bias, High variance
				\item Optimal: Balance between the two
			\end{itemize}
		\end{itemize}
		
		\vspace{0.05cm}
		\textcolor{myred}{\textbf{WHY OTHER OPTIONS ARE WRONG:}}
		\begin{itemize}
			\item \textcolor{myred}{A:} This is accuracy-interpretability tradeoff
			\item \textcolor{myred}{C:} This is speed-accuracy tradeoff
			\item \textcolor{myred}{D:} This is curse of dimensionality
		\end{itemize}
		
		\textcolor{myred}{\textbf{EXAM TRAP:}} Bias = error from assumptions; Variance = error from sensitivity to training data!
	\end{frame}
	
	\begin{frame}{Q47: Bias Definition - Tutorial}
		\fontsize{6.5}{7.5}\selectfont
		\textbf{What does "bias" refer to in the bias-variance tradeoff?}
		
		\vspace{0.05cm}
		\textcolor{myblue}{\textbf{ANSWER: A}} - Error introduced by the model's assumptions and limitations.
		
		\vspace{0.05cm}
		\textbf{DETAILED EXPLANATION:}
		\begin{itemize}
			\item \textbf{Definition:} Error due to model's simplifying assumptions
			\begin{itemize}
				\item Assumption of linearity when relationship is nonlinear
				\item Assumption of independence when data has structure
				\item Limited capacity of the model
			\end{itemize}
			\item \textbf{Mathematical definition:}
			\begin{itemize}
				\item $\text{Bias} = E[\hat{f}(X)] - f(X)$
				\item Expected difference between predicted and true function
				\item Over multiple training sets
			\end{itemize}
			\item \textbf{High bias characteristics:}
			\begin{itemize}
				\item Underfitting
				\item Poor performance on training data
				\item Poor performance on test data
				\item Simple model (linear, shallow)
			\end{itemize}
			\item \textbf{Examples of high bias models:}
			\begin{itemize}
				\item Linear regression for nonlinear data
				\item Logistic regression for complex classification
				\item Small decision trees
			\end{itemize}
		\end{itemize}
		
		\vspace{0.05cm}
		\textcolor{myred}{\textbf{WHY OTHER OPTIONS ARE WRONG:}}
		\begin{itemize}
			\item \textcolor{myred}{B:} This describes variance (sensitivity to training data)
			\item \textcolor{myred}{C:} This describes irreducible error
			\item \textcolor{myred}{D:} This describes optimization error
		\end{itemize}
		
		\textcolor{myred}{\textbf{EXAM TRAP:}} Bias = error from model assumptions (underfitting)!
	\end{frame}
	
	\begin{frame}{Q48: Variance Definition - Tutorial}
		\fontsize{6.5}{7.5}\selectfont
		\textbf{What does "variance" refer to in the bias-variance tradeoff?}
		
		\vspace{0.05cm}
		\textcolor{myblue}{\textbf{ANSWER: B}} - Error introduced by model's sensitivity to fluctuations in training data.
		
		\vspace{0.05cm}
		\textbf{DETAILED EXPLANATION:}
		\begin{itemize}
			\item \textbf{Definition:} Error due to sensitivity to training data
			\begin{itemize}
				\item Different training sets produce different models
				\item High variance models change a lot with data
				\item Captures noise as well as signal
			\end{itemize}
			\item \textbf{Mathematical definition:}
			\begin{itemize}
				\item $\text{Variance} = E[(\hat{f}(X) - E[\hat{f}(X)])^2]$
				\item Variability of predictions around expected prediction
				\item Over multiple training sets
			\end{itemize}
			\item \textbf{High variance characteristics:}
			\begin{itemize}
				\item Overfitting
				\item Excellent performance on training data
				\item Poor performance on test data
				\item Complex model (deep, many parameters)
			\end{itemize}
			\item \textbf{Examples of high variance models:}
			\begin{itemize}
				\item High-degree polynomial with too many terms
				\item Deep neural networks on small datasets
				\item Random forests with deep trees (without regularization)
			\end{itemize}
		\end{itemize}
		
		\vspace{0.05cm}
		\textcolor{myred}{\textbf{WHY OTHER OPTIONS ARE WRONG:}}
		\begin{itemize}
			\item \textcolor{myred}{A:} This describes bias (model assumptions)
			\item \textcolor{myred}{C:} This describes irreducible error
			\item \textcolor{myred}{D:} This describes a different type of error
		\end{itemize}
		
		\textcolor{myred}{\textbf{EXAM TRAP:}} Variance = error from sensitivity to training data (overfitting)!
	\end{frame}
	
	\begin{frame}{Q49: High Bias Characteristics - Tutorial}
		\fontsize{6.5}{7.5}\selectfont
		\textbf{Which is a characteristic of a model with high bias?}
		
		\vspace{0.05cm}
		\textcolor{myblue}{\textbf{ANSWER: B}} - Performs poorly on both training and test data (underfitting).
		
		\vspace{0.05cm}
		\textbf{DETAILED EXPLANATION:}
		\begin{itemize}
			\item \textbf{High bias pattern:}
			\begin{itemize}
				\item Training error: HIGH
				\item Test error: HIGH
				\item Both errors are similar
				\item Model is too simple
			\end{itemize}
			\item \textbf{Why this happens:}
			\begin{itemize}
				\item Model makes strong assumptions
				\item Cannot capture underlying patterns
				\item Underfits the data
			\end{itemize}
			\item \textbf{Diagnosing high bias:}
			\begin{itemize}
				\item Low training accuracy (for classification)
				\item High training loss
				\item No improvement with more data
				\item Model predictions are overly smooth
			\end{itemize}
			\item \textbf{Examples of high bias:}
			\begin{itemize}
				\item Linear regression on quadratic data
				\item Logistic regression with too few features
				\item Shallow decision trees
				\item Underparameterized neural networks
			\end{itemize}
			\item \textbf{Fixing high bias:}
			\begin{itemize}
				\item Increase model complexity
				\item Add more features
				\item Use more expressive model
				\item Reduce regularization
			\end{itemize}
		\end{itemize}
		
		\vspace{0.05cm}
		\textcolor{myred}{\textbf{WHY OTHER OPTIONS ARE WRONG:}}
		\begin{itemize}
			\item \textcolor{myred}{A:} This is low bias (fits training data well)
			\item \textcolor{myred}{C:} This is high variance (overfitting)
			\item \textcolor{myred}{D:} High bias is associated with low variance
		\end{itemize}
		
		\textcolor{myred}{\textbf{EXAM TRAP:}} High bias = poor performance on both training and test!
	\end{frame}
	
	\begin{frame}{Q50: High Variance Characteristics - Tutorial}
		\fontsize{6.5}{7.5}\selectfont
		\textbf{Which is a characteristic of a model with high variance?}
		
		\vspace{0.05cm}
		\textcolor{myblue}{\textbf{ANSWER: B}} - Performs very well on training data but poorly on test data (overfitting).
		
		\vspace{0.05cm}
		\textbf{DETAILED EXPLANATION:}
		\begin{itemize}
			\item \textbf{High variance pattern:}
			\begin{itemize}
				\item Training error: LOW
				\item Test error: HIGH
				\item Large gap between errors
				\item Model is too complex
			\end{itemize}
			\item \textbf{Why this happens:}
			\begin{itemize}
				\item Model memorizes training data
				\item Learns noise in addition to signal
				\item Overfits to specific patterns
			\end{itemize}
			\item \textbf{Diagnosing high variance:}
			\begin{itemize}
				\item High training accuracy, low test accuracy
				\item Model performance drops significantly on new data
				\item Large coefficients (for linear models)
				\item Complex decision boundaries
			\end{itemize}
			\item \textbf{Examples of high variance:}
			\begin{itemize}
				\item High-degree polynomial on small dataset
				\item Deep neural networks with too many parameters
				\item Decision trees with unlimited depth
				\item Any model with too many parameters for data size
			\end{itemize}
			\item \textbf{Fixing high variance:}
			\begin{itemize}
				\item Add regularization (L1, L2, dropout)
				\item Get more training data
				\item Reduce model complexity
				\item Use cross-validation
			\end{itemize}
		\end{itemize}
		
		\vspace{0.05cm}
		\textcolor{myred}{\textbf{WHY OTHER OPTIONS ARE WRONG:}}
		\begin{itemize}
			\item \textcolor{myred}{A:} This is high bias (underfitting)
			\item \textcolor{myred}{C:} High variance is associated with high complexity
			\item \textcolor{myred}{D:} High variance is associated with low bias
		\end{itemize}
		
		\textcolor{myred}{\textbf{EXAM TRAP:}} High variance = good training, poor test (overfitting)!
	\end{frame}
	
	\begin{frame}{Q51: Underfitting Definition - Tutorial}
		\fontsize{6.5}{7.5}\selectfont
		\textbf{What is underfitting in supervised learning?}
		
		\vspace{0.05cm}
		\textcolor{myblue}{\textbf{ANSWER: B}} - Model is too simple to capture underlying patterns in data.
		
		\vspace{0.05cm}
		\textbf{DETAILED EXPLANATION:}
		\begin{itemize}
			\item \textbf{Definition:} Model fails to capture data's underlying structure
			\begin{itemize}
				\item Too simple for the task
				\item Cannot represent the true function
				\item Makes overly strong assumptions
			\end{itemize}
			\item \textbf{Signs of underfitting:}
			\begin{itemize}
				\item High training error
				\item High test error
				\item No significant gap between training and test error
				\item Predictions are too simplistic
			\end{itemize}
			\item \textbf{Causes:}
			\begin{itemize}
				\item Model is too simple (linear for nonlinear data)
				\item Too few features
				\item Too much regularization (too large $\lambda$)
				\item Insufficient training (early stopping too early)
			\end{itemize}
			\item \textbf{Fixing underfitting:}
			\begin{itemize}
				\item Increase model complexity (more parameters)
				\item Add more features
				\item Decrease regularization
				\item Train for more iterations
			\end{itemize}
			\item \textbf{Analogy:} Trying to fit a sine wave with a straight line
		\end{itemize}
		
		\vspace{0.05cm}
		\textcolor{myred}{\textbf{WHY OTHER OPTIONS ARE WRONG:}}
		\begin{itemize}
			\item \textcolor{myred}{A:} This describes overfitting
			\item \textcolor{myred}{C:} This describes overfitting
			\item \textcolor{myred}{D:} Too many parameters leads to overfitting
		\end{itemize}
		
		\textcolor{myred}{\textbf{EXAM TRAP:}} Underfitting = model too simple = poor training and test performance!
	\end{frame}
	
	\begin{frame}{Q52: Overfitting Definition - Tutorial}
		\fontsize{6.5}{7.5}\selectfont
		\textbf{What is overfitting in supervised learning?}
		
		\vspace{0.05cm}
		\textcolor{myblue}{\textbf{ANSWER: B}} - Model is too complex and learns the noise in the training data.
		
		\vspace{0.05cm}
		\textbf{DETAILED EXPLANATION:}
		\begin{itemize}
			\item \textbf{Definition:} Model fits training data too well, including noise
			\begin{itemize}
				\item Memorizes training examples
				\item Cannot generalize to new data
				\item Captures random fluctuations
			\end{itemize}
			\item \textbf{Signs of overfitting:}
			\begin{itemize}
				\item Very low training error
				\item High test error
				\item Large gap between training and test error
				\item Model has many parameters relative to data
			\end{itemize}
			\item \textbf{Causes:}
			\begin{itemize}
				\item Model is too complex
				\item Too many features
				\item Too little regularization (too small $\lambda$)
				\item Too little training data
			\end{itemize}
			\item \textbf{Fixing overfitting:}
			\begin{itemize}
				\item Add regularization (L1, L2, dropout)
				\item Get more training data
				\item Reduce model complexity
				\item Use cross-validation
				\item Early stopping
			\end{itemize}
			\item \textbf{Analogy:} Memorizing answers instead of learning concepts
		\end{itemize}
		
		\vspace{0.05cm}
		\textcolor{myred}{\textbf{WHY OTHER OPTIONS ARE WRONG:}}
		\begin{itemize}
			\item \textcolor{myred}{A:} This describes underfitting
			\item \textcolor{myred}{C:} This describes underfitting
			\item \textcolor{myred}{D:} Too few parameters leads to underfitting
		\end{itemize}
		
		\textcolor{myred}{\textbf{EXAM TRAP:}} Overfitting = model too complex = good training, poor test!
	\end{frame}
	
	\begin{frame}{Q53: Causes of Overfitting - Tutorial}
		\fontsize{6.5}{7.5}\selectfont
		\textbf{Which is a common cause of overfitting?}
		
		\vspace{0.05cm}
		\textcolor{myblue}{\textbf{ANSWER: B}} - Too many parameters relative to number of training samples.
		
		\vspace{0.05cm}
		\textbf{DETAILED EXPLANATION:}
		\begin{itemize}
			\item \textbf{The core issue:} $p \gg N$ (more parameters than samples)
			\begin{itemize}
				\item More parameters than constraints
				\item Infinite solutions that fit data perfectly
				\item Model can memorize individual examples
			\end{itemize}
			\item \textbf{Example scenarios:}
			\begin{itemize}
				\item 100 parameters with 50 samples
				\item High-degree polynomial (degree 20) with 30 points
				\item Deep neural network on small dataset
			\end{itemize}
			\item \textbf{Other contributing factors:}
			\begin{itemize}
				\item Too many features (high-dimensional)
				\item Little or no regularization
				\item Training for too many iterations
				\item Data with high noise
			\end{itemize}
			\item \textbf{Prevention strategies:}
			\begin{itemize}
				\item Use cross-validation
				\item Apply regularization
				\item Reduce model complexity
				\item Get more training data
				\item Use feature selection
			\end{itemize}
		\end{itemize}
		
		\vspace{0.05cm}
		\textcolor{myred}{\textbf{WHY OTHER OPTIONS ARE WRONG:}}
		\begin{itemize}
			\item \textcolor{myred}{A:} Too few features leads to underfitting
			\item \textcolor{myred}{C:} Too much regularization leads to underfitting
			\item \textcolor{myred}{D:} Too simple a model leads to underfitting
		\end{itemize}
		
		\textcolor{myred}{\textbf{EXAM TRAP:}} Overfitting happens when too many parameters for too few samples!
	\end{frame}
	
	\begin{frame}{Q54: Causes of Underfitting - Tutorial}
		\fontsize{6.5}{7.5}\selectfont
		\textbf{Which is a common cause of underfitting?}
		
		\vspace{0.05cm}
		\textcolor{myblue}{\textbf{ANSWER: B}} - Insufficient model complexity (e.g., linear for nonlinear data).
		
		\vspace{0.05cm}
		\textbf{DETAILED EXPLANATION:}
		\begin{itemize}
			\item \textbf{The core issue:} Model cannot represent true function
			\begin{itemize}
				\item Model class is too restrictive
				\item Cannot capture the data's patterns
				\item Relationship is more complex than model allows
			\end{itemize}
			\item \textbf{Common scenarios:}
			\begin{itemize}
				\item Linear regression for quadratic relationship
				\item Logistic regression for complex decision boundary
				\item Shallow network for deep patterns
				\item Too few features
			\end{itemize}
			\item \textbf{Other contributing factors:}
			\begin{itemize}
				\item Too much regularization ($\lambda$ too large)
				\item Stopping training too early
				\item Poor feature engineering
				\item Too simple model architecture
			\end{itemize}
			\item \textbf{Fixing underfitting:}
			\begin{itemize}
				\item Choose more expressive model
				\item Add more features
				\item Decrease regularization
				\item Train for more iterations
				\item Use better feature engineering
			\end{itemize}
		\end{itemize}
		
		\vspace{0.05cm}
		\textcolor{myred}{\textbf{WHY OTHER OPTIONS ARE WRONG:}}
		\begin{itemize}
			\item \textcolor{myred}{A:} Too many parameters leads to overfitting
			\item \textcolor{myred}{C:} Too little regularization leads to overfitting
			\item \textcolor{myred}{D:} More training data generally helps
		\end{itemize}
		
		\textcolor{myred}{\textbf{EXAM TRAP:}} Underfitting happens when model is too simple for the data!
	\end{frame}
	
	\begin{frame}{Q55: Bias-Variance Tradeoff - Tutorial}
		\fontsize{6.5}{7.5}\selectfont
		\textbf{As model complexity increases, what happens to bias and variance?}
		
		\vspace{0.05cm}
		\textcolor{myblue}{\textbf{ANSWER: B}} - Bias decreases and variance increases.
		
		\vspace{0.05cm}
		\textbf{DETAILED EXPLANATION:}
		\begin{itemize}
			\item \textbf{Relationship:}
			\begin{itemize}
				\item \textbf{Bias:} $\downarrow$ as complexity increases
				\item \textbf{Variance:} $\uparrow$ as complexity increases
				\item \textbf{Total error:} U-shaped curve
			\end{itemize}
			\item \textbf{Why bias decreases:}
			\begin{itemize}
				\item More flexible model captures more patterns
				\item Can represent more complex functions
				\item Assumptions are less restrictive
			\end{itemize}
			\item \textbf{Why variance increases:}
			\begin{itemize}
				\item Model becomes more sensitive to data
				\item More parameters mean more degrees of freedom
				\item Can fit noise as well as signal
			\end{itemize}
			\item \textbf{The optimal point:}
			\begin{itemize}
				\item Minimum of bias$^2$ + variance
				\item Balance between underfitting and overfitting
				\item Tradeoff is fundamental
			\end{itemize}
		\end{itemize}
		
		\vspace{0.05cm}
		\textcolor{myred}{\textbf{WHY OTHER OPTIONS ARE WRONG:}}
		\begin{itemize}
			\item \textcolor{myred}{A:} Both decrease is incorrect
			\item \textcolor{myred}{C:} This is the opposite direction
			\item \textcolor{myred}{D:} Both increase is incorrect
		\end{itemize}
		
		\textcolor{myred}{\textbf{EXAM TRAP:}} More complexity $\to$ lower bias, higher variance!
	\end{frame}
	
	\begin{frame}{Q56: Optimal Complexity - Tutorial}
		\fontsize{6.5}{7.5}\selectfont
		\textbf{Where is the optimal model complexity located?}
		
		\vspace{0.05cm}
		\textcolor{myblue}{\textbf{ANSWER: B}} - Where total error (bias$^2$ + variance + irreducible) is minimized.
		
		\vspace{0.05cm}
		\textbf{DETAILED EXPLANATION:}
		\begin{itemize}
			\item \textbf{Total error decomposition:}
			\begin{itemize}
				\item Total Error = Bias$^2$ + Variance + Irreducible Error
				\item Bias$^2$: decreases with complexity
				\item Variance: increases with complexity
				\item Irreducible: constant
			\end{itemize}
			\item \textbf{Shape of error curve:}
			\begin{itemize}
				\item Low complexity: High bias, Low variance $\to$ High total error
				\item High complexity: Low bias, High variance $\to$ High total error
				\item Optimal: Sweet spot in between
			\end{itemize}
			\item \textbf{Finding optimal complexity:}
			\begin{itemize}
				\item Use cross-validation
				\item Track validation error
				\item Choose model with lowest validation error
			\end{itemize}
			\item \textbf{Practical considerations:}
			\begin{itemize}
				\item Optimal complexity depends on data size
				\item More data allows for more complex models
				\item Regularization helps shift the optimal point
			\end{itemize}
		\end{itemize}
		
		\vspace{0.05cm}
		\textcolor{myred}{\textbf{WHY OTHER OPTIONS ARE WRONG:}}
		\begin{itemize}
			\item \textcolor{myred}{A:} Highest variance = overfitting (too complex)
			\item \textcolor{myred}{C:} Lowest bias = overfitting (too complex)
			\item \textcolor{myred}{D:} Highest bias = underfitting (too simple)
		\end{itemize}
		
		\textcolor{myred}{\textbf{EXAM TRAP:}} Optimal complexity = minimum total error (bias$^2$ + variance)!
	\end{frame}
	
	\begin{frame}{Q57: Bias-Variance Example - Tutorial}
		\fontsize{6.5}{7.5}\selectfont
		\textbf{Simple linear regression for nonlinear relationship would likely have:}
		
		\vspace{0.05cm}
		\textcolor{myblue}{\textbf{ANSWER: A}} - High bias, low variance (underfitting).
		
		\vspace{0.05cm}
		\textbf{DETAILED EXPLANATION:}
		\begin{itemize}
			\item \textbf{Why high bias:}
			\begin{itemize}
				\item Linear model assumes linear relationship
				\item True relationship is nonlinear
				\item Model cannot capture curvature
				\item Systematically wrong predictions
			\end{itemize}
			\item \textbf{Why low variance:}
			\begin{itemize}
				\item Linear model is very simple
				\item Few parameters (slope and intercept)
				\item Not sensitive to training data
				\item Stable predictions across datasets
			\end{itemize}
			\item \textbf{Result:} Underfitting
			\begin{itemize}
				\item High training error
				\item High test error
				\item Model is too simple
			\end{itemize}
			\item \textbf{Fix:} Use polynomial regression or other nonlinear model
		\end{itemize}
		
		\vspace{0.05cm}
		\textcolor{myred}{\textbf{WHY OTHER OPTIONS ARE WRONG:}}
		\begin{itemize}
			\item \textcolor{myred}{B:} This describes complex model (overfitting)
			\item \textcolor{myred}{C:} Ideal but unrealistic for mismatched model
			\item \textcolor{myred}{D:} Unlikely combination
		\end{itemize}
		
		\textcolor{myred}{\textbf{EXAM TRAP:}} Too simple $\to$ high bias, low variance (underfitting)!
	\end{frame}
	
	\begin{frame}{Q58: High-Degree Polynomial - Tutorial}
		\fontsize{6.5}{7.5}\selectfont
		\textbf{High-degree polynomial on small dataset would likely have:}
		
		\vspace{0.05cm}
		\textcolor{myblue}{\textbf{ANSWER: B}} - Low bias, high variance (overfitting).
		
		\vspace{0.05cm}
		\textbf{DETAILED EXPLANATION:}
		\begin{itemize}
			\item \textbf{Why low bias:}
			\begin{itemize}
				\item High-degree polynomial is very flexible
				\item Can fit any finite set of points
				\item Captures all patterns in training data
				\item Very expressive model
			\end{itemize}
			\item \textbf{Why high variance:}
			\begin{itemize}
				\item Many parameters (degree + 1)
				\item Very sensitive to training data
				\item Different dataset $\to$ very different fit
				\item Oscillates between data points
			\end{itemize}
			\item \textbf{Result:} Overfitting
			\begin{itemize}
				\item Very low training error
				\item High test error
				\item Model is too complex
			\end{itemize}
			\item \textbf{Fix:} Use regularization or simpler model
		\end{itemize}
		
		\vspace{0.05cm}
		\textcolor{myred}{\textbf{WHY OTHER OPTIONS ARE WRONG:}}
		\begin{itemize}
			\item \textcolor{myred}{A:} This describes simple model (underfitting)
			\item \textcolor{myred}{C:} Ideal but unrealistic for complex model on small dataset
			\item \textcolor{myred}{D:} Unlikely combination
		\end{itemize}
		
		\textcolor{myred}{\textbf{EXAM TRAP:}} Too complex $\to$ low bias, high variance (overfitting)!
	\end{frame}
	
	\begin{frame}{Q59: Irreducible Error - Tutorial}
		\fontsize{6.5}{7.5}\selectfont
		\textbf{What is irreducible error in the bias-variance decomposition?}
		
		\vspace{0.05cm}
		\textcolor{myblue}{\textbf{ANSWER: C}} - Error due to noise in data that cannot be reduced by any model.
		
		\vspace{0.05cm}
		\textbf{DETAILED EXPLANATION:}
		\begin{itemize}
			\item \textbf{Definition:} Inherent noise in the data
			\begin{itemize}
				\item Measurement errors
				\item Unobserved variables
				\item Inherent randomness
				\item Natural variation
			\end{itemize}
			\item \textbf{Mathematical notation:}
			\begin{itemize}
				\item $Y = f(X) + \varepsilon$
				\item $\varepsilon \sim \text{some distribution}$
				\item Var$(\varepsilon) = \sigma^2$ = irreducible error
			\end{itemize}
			\item \textbf{Why irreducible:}
			\begin{itemize}
				\item No model can eliminate this error
				\item Even if we knew true $f(X)$, still have $\varepsilon$
				\item Best possible model still has error $\sigma^2$
			\end{itemize}
			\item \textbf{Practical implications:}
			\begin{itemize}
				\item There is always a limit to predictive accuracy
				\item Cannot achieve zero error
				\item Need to accept some level of error
			\end{itemize}
		\end{itemize}
		
		\vspace{0.05cm}
		\textcolor{myred}{\textbf{WHY OTHER OPTIONS ARE WRONG:}}
		\begin{itemize}
			\item \textcolor{myred}{A:} This describes bias (reducible error)
			\item \textcolor{myred}{B:} This describes variance (reducible error)
			\item \textcolor{myred}{D:} Loss function choice is a modeling decision
		\end{itemize}
		
		\textcolor{myred}{\textbf{EXAM TRAP:}} Irreducible error = noise in data - cannot be eliminated!
	\end{frame}
	
	\begin{frame}{Q60: Bias-Variance Summary - Tutorial}
		\fontsize{6.5}{7.5}\selectfont
		\textbf{Which statement about the bias-variance tradeoff is TRUE?}
		
		\vspace{0.05cm}
		\textcolor{myblue}{\textbf{ANSWER: C}} - The optimal model balances bias and variance to minimize total error.
		
		\vspace{0.05cm}
		\textbf{DETAILED EXPLANATION:}
		\begin{itemize}
			\item \textbf{Key concept:} Balance is the goal
			\begin{itemize}
				\item Not too simple (high bias)
				\item Not too complex (high variance)
				\item Find the sweet spot
			\end{itemize}
			\item \textbf{Why other options are false:}
			\begin{itemize}
				\item \textcolor{myred}{A:} Simple models have high bias, low variance (opposite)
				\item \textcolor{myred}{B:} Complex models have low bias, high variance (opposite)
				\item \textcolor{myred}{D:} Bias and variance are related through complexity
			\end{itemize}
			\item \textbf{How to balance:}
			\begin{itemize}
				\item Use cross-validation to select model
				\item Use regularization to control complexity
				\item Monitor training and validation error
			\end{itemize}
			\item \textbf{Real-world implications:}
			\begin{itemize}
				\item No single best model for all problems
				\item Tradeoff is fundamental
				\item Choice depends on problem and data
			\end{itemize}
		\end{itemize}
		
		\vspace{0.05cm}
		\textcolor{myred}{\textbf{WHY OTHER OPTIONS ARE WRONG:}}
		\begin{itemize}
			\item \textcolor{myred}{A:} Simple models have high bias, low variance
			\item \textcolor{myred}{B:} Complex models have low bias, high variance
			\item \textcolor{myred}{D:} Bias and variance are related through model complexity
		\end{itemize}
		
		\textcolor{myred}{\textbf{EXAM TRAP:}} Optimal model balances bias and variance!
	\end{frame}
	
	% ============================================================
	% SECTION 5: REGULARIZATION - QUESTIONS 61-75
	% ============================================================
	
	\begin{frame}{Q61: Regularization Definition - Tutorial}
		\fontsize{6.5}{7.5}\selectfont
		\textbf{What is regularization in supervised learning?}
		
		\vspace{0.05cm}
		\textcolor{myblue}{\textbf{ANSWER: B}} - Technique to reduce overfitting by adding a penalty term to the loss function.
		
		\vspace{0.05cm}
		\textbf{DETAILED EXPLANATION:}
		\begin{itemize}
			\item \textbf{Definition:} Adding penalty to prevent overfitting
			\begin{itemize}
				\item $\tilde{L}(\theta) = L(\theta) + \lambda \cdot \text{Penalty}(\theta)$
				\item $L(\theta)$: original loss
				\item $\lambda$: regularization strength
				\item Penalty: encourages simpler models
			\end{itemize}
			\item \textbf{Why it works:}
			\begin{itemize}
				\item Penalizes large parameters
				\item Reduces model complexity
				\item Improves generalization
				\item Controls the bias-variance tradeoff
			\end{itemize}
			\item \textbf{Types of regularization:}
			\begin{itemize}
				\item L1 (Lasso): $\lambda \sum |w_i|$
				\item L2 (Ridge): $\lambda \sum w_i^2$
				\item Elastic Net: combination
				\item Dropout (neural networks)
				\item Early stopping
			\end{itemize}
		\end{itemize}
		
		\vspace{0.05cm}
		\textcolor{myred}{\textbf{WHY OTHER OPTIONS ARE WRONG:}}
		\begin{itemize}
			\item \textcolor{myred}{A:} This is the opposite (increasing complexity)
			\item \textcolor{myred}{C:} This is data augmentation
			\item \textcolor{myred}{D:} This is feature selection (related but distinct)
		\end{itemize}
		
		\textcolor{myred}{\textbf{EXAM TRAP:}} Regularization adds a penalty to the loss function!
	\end{frame}
	
	\begin{frame}{Q62: L1 Regularization (Lasso) - Tutorial}
		\fontsize{6.5}{7.5}\selectfont
		\textbf{What is L1 regularization (Lasso) and its key property?}
		
		\vspace{0.05cm}
		\textcolor{myblue}{\textbf{ANSWER: B}} - Adds $\lambda \sum |w_i|$ to loss; can set weights to exactly zero (feature selection).
		
		\vspace{0.05cm}
		\textbf{DETAILED EXPLANATION:}
		\begin{itemize}
			\item \textbf{Lasso objective:}
			\begin{itemize}
				\item $\tilde{L}(\theta) = L(\theta) + \lambda \sum_{j=1}^p |w_j|$
				\item $L(\theta)$: original loss (e.g., MSE)
				\item $\lambda$: regularization strength
			\end{itemize}
			\item \textbf{Key property: Feature selection}
			\begin{itemize}
				\item Can set weights to exactly zero
				\item Effectively removes features
				\item Sparse solutions
				\item Identifies important features
			\end{itemize}
			\item \textbf{Why it works:}
			\begin{itemize}
				\item L1 penalty is non-differentiable at zero
				\item Creates "corners" where weights become zero
				\item Geometry: diamond-shaped constraint region
			\end{itemize}
			\item \textbf{When to use:}
			\begin{itemize}
				\item Many features, few important ones
				\item Want interpretable model
				\item Feature selection needed
			\end{itemize}
		\end{itemize}
		
		\vspace{0.05cm}
		\textcolor{myred}{\textbf{WHY OTHER OPTIONS ARE WRONG:}}
		\begin{itemize}
			\item \textcolor{myred}{A:} This is L2 regularization (Ridge)
			\item \textcolor{myred}{C:} Linear penalty is not standard
			\item \textcolor{myred}{D:} Cubic penalty is not standard regularization
		\end{itemize}
		
		\textcolor{myred}{\textbf{EXAM TRAP:}} L1 regularization can set weights to zero (feature selection)!
	\end{frame}
	
	\begin{frame}{Q63: L2 Regularization (Ridge) - Tutorial}
		\fontsize{6.5}{7.5}\selectfont
		\textbf{What is L2 regularization (Ridge) and its key property?}
		
		\vspace{0.05cm}
		\textcolor{myblue}{\textbf{ANSWER: B}} - Adds $\lambda \sum w_i^2$ to loss; shrinks weights toward zero but does not set them to zero.
		
		\vspace{0.05cm}
		\textbf{DETAILED EXPLANATION:}
		\begin{itemize}
			\item \textbf{Ridge objective:}
			\begin{itemize}
				\item $\tilde{L}(\theta) = L(\theta) + \lambda \sum_{j=1}^p w_j^2$
				\item $L(\theta)$: original loss (e.g., MSE)
				\item $\lambda$: regularization strength
			\end{itemize}
			\item \textbf{Key property: Weight shrinkage}
			\begin{itemize}
				\item Shrinks weights toward zero
				\item Never exactly zero (unless $\lambda = \infty$)
				\item Dense solutions
				\item All features retained
			\end{itemize}
			\item \textbf{Why it works:}
			\begin{itemize}
				\item L2 penalty is differentiable everywhere
				\item Geometry: circular constraint region
				\item Smooth shrinkage
			\end{itemize}
			\item \textbf{When to use:}
			\begin{itemize}
				\item Many features, all potentially important
				\item Multicollinearity present
				\item Want to keep all features
				\item Better predictive performance than L1 sometimes
			\end{itemize}
		\end{itemize}
		
		\vspace{0.05cm}
		\textcolor{myred}{\textbf{WHY OTHER OPTIONS ARE WRONG:}}
		\begin{itemize}
			\item \textcolor{myred}{A:} This is L1 regularization (Lasso)
			\item \textcolor{myred}{C:} Linear penalty is not standard
			\item \textcolor{myred}{D:} Cubic penalty is not standard regularization
		\end{itemize}
		
		\textcolor{myred}{\textbf{EXAM TRAP:}} L2 regularization shrinks weights but does not set them to zero!
	\end{frame}
	
	\begin{frame}{Q64: Elastic Net - Tutorial}
		\fontsize{6.5}{7.5}\selectfont
		\textbf{What is Elastic Net regularization?}
		
		\vspace{0.05cm}
		\textcolor{myblue}{\textbf{ANSWER: A}} - Combination of L1 and L2: $\lambda_1 \sum |w_i| + \lambda_2 \sum w_i^2$
		
		\vspace{0.05cm}
		\textbf{DETAILED EXPLANATION:}
		\begin{itemize}
			\item \textbf{Elastic Net objective:}
			\begin{itemize}
				\item $\tilde{L}(\theta) = L(\theta) + \lambda_1 \sum |w_i| + \lambda_2 \sum w_i^2$
				\item Combines L1 and L2 penalties
				\item Two hyperparameters: $\lambda_1$ and $\lambda_2$
			\end{itemize}
			\item \textbf{Advantages:}
			\begin{itemize}
				\item \textbf{Feature selection} (from L1)
				\item \textbf{Group selection}: Correlated features enter together
				\item \textbf{Better than Lasso} when features are correlated
				\item \textbf{Better than Ridge} when only some features important
			\end{itemize}
			\item \textbf{Why it's needed:}
			\begin{itemize}
				\item Lasso struggles with correlated features (selects one randomly)
				\item Ridge keeps all features
				\item Elastic Net: correlation-aware feature selection
			\end{itemize}
			\item \textbf{Common parameterization:}
			\begin{itemize}
				\item $\alpha$: L1 ratio ($\alpha = 1$ = Lasso, $\alpha = 0$ = Ridge)
				\item $\lambda$: overall regularization strength
				\item Penalty: $\lambda(\alpha \sum |w_i| + (1-\alpha)\sum w_i^2)$
			\end{itemize}
		\end{itemize}
		
		\vspace{0.05cm}
		\textcolor{myred}{\textbf{WHY OTHER OPTIONS ARE WRONG:}}
		\begin{itemize}
			\item \textcolor{myred}{B:} This is Lasso (L1 only)
			\item \textcolor{myred}{C:} This is Ridge (L2 only)
			\item \textcolor{myred}{D:} This is no regularization
		\end{itemize}
		
		\textcolor{myred}{\textbf{EXAM TRAP:}} Elastic Net = L1 + L2 regularization!
	\end{frame}
	
	\begin{frame}{Q65: Regularization Parameter $\lambda$ - Tutorial}
		\fontsize{6.5}{7.5}\selectfont
		\textbf{What happens as the regularization parameter $\lambda$ increases?}
		
		\vspace{0.05cm}
		\textcolor{myblue}{\textbf{ANSWER: B}} - The model becomes simpler (weights shrink more toward zero).
		
		\vspace{0.05cm}
		\textbf{DETAILED EXPLANATION:}
		\begin{itemize}
			\item \textbf{Effect of increasing $\lambda$:}
			\begin{itemize}
				\item Stronger penalty on large weights
				\item Weights shrink closer to zero
				\item Model becomes more constrained
				\item Bias increases, variance decreases
			\end{itemize}
			\item \textbf{Extreme cases:}
			\begin{itemize}
				\item $\lambda = 0$: No regularization (standard OLS)
				\item $\lambda \to \infty$: All weights $\to 0$ (underfitting)
			\end{itemize}
			\item \textbf{Choosing $\lambda$:}
			\begin{itemize}
				\item Use cross-validation
				\item Balance between underfitting and overfitting
				\item Typically choose $\lambda$ that minimizes validation error
			\end{itemize}
			\item \textbf{Model complexity relationship:}
			\begin{itemize}
				\item $\lambda$ controls model complexity
				\item Larger $\lambda$ $\to$ simpler model
				\item Smaller $\lambda$ $\to$ more complex model
			\end{itemize}
		\end{itemize}
		
		\vspace{0.05cm}
		\textcolor{myred}{\textbf{WHY OTHER OPTIONS ARE WRONG:}}
		\begin{itemize}
			\item \textcolor{myred}{A:} Larger $\lambda$ means more regularization = simpler model
			\item \textcolor{myred}{C:} Larger $\lambda$ reduces overfitting
			\item \textcolor{myred}{D:} $\lambda$ definitely changes the model
		\end{itemize}
		
		\textcolor{myred}{\textbf{EXAM TRAP:}} Larger $\lambda$ $\to$ simpler model (more regularization)!
	\end{frame}
	
	\begin{frame}{Q66: Regularization Overfitting Prevention - Tutorial}
		\fontsize{6.5}{7.5}\selectfont
		\textbf{How does regularization prevent overfitting?}
		
		\vspace{0.05cm}
		\textcolor{myblue}{\textbf{ANSWER: B}} - By penalizing large parameter values, reducing model complexity and sensitivity to training data.
		
		\vspace{0.05cm}
		\textbf{DETAILED EXPLANATION:}
		\begin{itemize}
			\item \textbf{Mechanism:}
			\begin{itemize}
				\item Large weights $\to$ complex decision boundaries
				\item Large weights $\to$ high sensitivity to inputs
				\item Penalty discourages large weights
				\item Result: smoother, simpler models
			\end{itemize}
			\item \textbf{Why this prevents overfitting:}
			\begin{itemize}
				\item Model cannot fit noise (requires large weights)
				\item Less sensitive to training data fluctuations
				\item Better generalization to new data
			\end{itemize}
			\item \textbf{Examples:}
			\begin{itemize}
				\item L2: $\lambda \sum w_i^2$ penalizes squared weights
				\item L1: $\lambda \sum |w_i|$ penalizes absolute weights
				\item Both discourage large weights
			\end{itemize}
			\item \textbf{Analogy:} Occam's Razor - prefer simpler explanations
		\end{itemize}
		
		\vspace{0.05cm}
		\textcolor{myred}{\textbf{WHY OTHER OPTIONS ARE WRONG:}}
		\begin{itemize}
			\item \textcolor{myred}{A:} Increasing features increases overfitting
			\item \textcolor{myred}{C:} This is data augmentation
			\item \textcolor{myred}{D:} More complex loss isn't regularization
		\end{itemize}
		
		\textcolor{myred}{\textbf{EXAM TRAP:}} Regularization = penalizes large weights to reduce overfitting!
	\end{frame}
	
	\begin{frame}{Q67: Dropout as Regularization - Tutorial}
		\fontsize{6.5}{7.5}\selectfont
		\textbf{Is dropout a form of regularization? If so, how does it work?}
		
		\vspace{0.05cm}
		\textcolor{myblue}{\textbf{ANSWER: B}} - Yes, randomly drops neurons during training, creating ensemble effect and preventing co-adaptation.
		
		\vspace{0.05cm}
		\textbf{DETAILED EXPLANATION:}
		\begin{itemize}
			\item \textbf{Definition:} Dropout is a regularization technique for neural networks
			\begin{itemize}
				\item Randomly sets a fraction of neurons to zero
				\item During each training iteration
				\item Typical dropout rate: 0.2 to 0.5
			\end{itemize}
			\item \textbf{How it works:}
			\begin{itemize}
				\item \textbf{During training:} Each neuron has probability $p$ of being dropped
				\item \textbf{During testing:} All neurons active, scaled by $(1-p)$
				\item Prevents co-adaptation of neurons
			\end{itemize}
			\item \textbf{Why it's effective:}
			\begin{itemize}
				\item Creates an ensemble effect (many subnetworks)
				\item Forces network to learn robust features
				\item Prevents over-reliance on specific neurons
				\item Acts as regularization (increases bias, decreases variance)
			\end{itemize}
			\item \textbf{Practical notes:}
			\begin{itemize}
				\item Used with large neural networks
				\item Particularly effective in deep networks
				\item Often combined with L2 regularization
			\end{itemize}
		\end{itemize}
		
		\vspace{0.05cm}
		\textcolor{myred}{\textbf{WHY OTHER OPTIONS ARE WRONG:}}
		\begin{itemize}
			\item \textcolor{myred}{A:} False - dropout IS regularization
			\item \textcolor{myred}{C:} False - dropout doesn't add penalty term
			\item \textcolor{myred}{D:} False - dropout affects both forward and backward passes
		\end{itemize}
		
		\textcolor{myred}{\textbf{EXAM TRAP:}} Dropout is a regularization technique for neural networks!
	\end{frame}
	
	\begin{frame}{Q68: Early Stopping as Regularization - Tutorial}
		\fontsize{6.5}{7.5}\selectfont
		\textbf{Is early stopping a form of regularization? If so, how does it work?}
		
		\vspace{0.05cm}
		\textcolor{myblue}{\textbf{ANSWER: B}} - Yes, stops training before convergence, limiting model complexity and preventing overfitting.
		
		\vspace{0.05cm}
		\textbf{DETAILED EXPLANATION:}
		\begin{itemize}
			\item \textbf{Definition:} Early stopping stops training early
			\begin{itemize}
				\item Monitor validation error
				\item Stop when validation error starts increasing
				\item Prevents overfitting
			\end{itemize}
			\item \textbf{How it works:}
			\begin{itemize}
				\item Track training and validation error
				\item Validation error typically: decreases then increases
				\item Stop at minimum validation error
				\item Use parameters from best validation performance
			\end{itemize}
			\item \textbf{Why it's regularization:}
			\begin{itemize}
				\item Limits the number of training iterations
				\item Equivalent to early stopping in optimization
				\item Prevents fitting noise
				\item Controls model complexity implicitly
			\end{itemize}
			\item \textbf{Practical implementation:}
			\begin{itemize}
				\item Patience parameter: wait for $k$ epochs before stopping
				\item Save best model checkpoint
				\item Return to best checkpoint after stopping
			\end{itemize}
		\end{itemize}
		
		\vspace{0.05cm}
		\textcolor{myred}{\textbf{WHY OTHER OPTIONS ARE WRONG:}}
		\begin{itemize}
			\item \textcolor{myred}{A:} False - early stopping IS regularization
			\item \textcolor{myred}{C:} False - early stopping doesn't add penalty term
			\item \textcolor{myred}{D:} False - early stopping affects model (prevents overfitting)
		\end{itemize}
		
		\textcolor{myred}{\textbf{EXAM TRAP:}} Early stopping is a regularization technique!
	\end{frame}
	
	\begin{frame}{Q69: Regularization Bias-Variance - Tutorial}
		\fontsize{6.5}{7.5}\selectfont
		\textbf{How does regularization affect the bias-variance tradeoff?}
		
		\vspace{0.05cm}
		\textcolor{myblue}{\textbf{ANSWER: C}} - Regularization increases bias and decreases variance.
		
		\vspace{0.05cm}
		\textbf{DETAILED EXPLANATION:}
		\begin{itemize}
			\item \textbf{Effect on bias:}
			\begin{itemize}
				\item Regularization makes model simpler
				\item Model may not fit training data perfectly
				\item Systematic error (bias) increases
			\end{itemize}
			\item \textbf{Effect on variance:}
			\begin{itemize}
				\item Regularization constrains weights
				\item Model is less sensitive to training data
				\item Variability (variance) decreases
			\end{itemize}
			\item \textbf{The tradeoff:}
			\begin{itemize}
				\item Regularization shifts towards simpler models
				\item Reduces overfitting (decreases variance)
				\item Increases underfitting risk (increases bias)
				\item Optimal $\lambda$ balances the two
			\end{itemize}
			\item \textbf{Visual representation:}
			\begin{itemize}
				\item No regularization: low bias, high variance
				\item Too much regularization: high bias, low variance
				\item Optimal: balanced
			\end{itemize}
		\end{itemize}
		
		\vspace{0.05cm}
		\textcolor{myred}{\textbf{WHY OTHER OPTIONS ARE WRONG:}}
		\begin{itemize}
			\item \textcolor{myred}{A:} Variance decreases, not increases
			\item \textcolor{myred}{B:} This is the opposite direction
			\item \textcolor{myred}{D:} Regularization definitely affects both
		\end{itemize}
		
		\textcolor{myred}{\textbf{EXAM TRAP:}} Regularization $\to$ increases bias, decreases variance!
	\end{frame}
	
	\begin{frame}{Q70: When to Use Regularization - Tutorial}
		\fontsize{6.5}{7.5}\selectfont
		\textbf{When should regularization be applied in supervised learning?}
		
		\vspace{0.05cm}
		\textcolor{myblue}{\textbf{ANSWER: B}} - When the model is overfitting (high variance).
		
		\vspace{0.05cm}
		\textbf{DETAILED EXPLANATION:}
		\begin{itemize}
			\item \textbf{Indications for regularization:}
			\begin{itemize}
				\item Training error much lower than validation error
				\item Model performs well on training, poorly on test
				\item Large coefficients (for linear models)
				\item Complex model with limited data
			\end{itemize}
			\item \textbf{Why it works:}
			\begin{itemize}
				\item Regularization reduces variance
				\item Prevents model from fitting noise
				\item Improves generalization
			\end{itemize}
			\item \textbf{When NOT to use regularization:}
			\begin{itemize}
				\item Model is underfitting (high bias)
				\item Training error is already high
				\item Data is large enough for model complexity
				\item Model is interpretable and simple
			\end{itemize}
			\item \textbf{Practical advice:}
			\begin{itemize}
				\item Always consider regularization for high-dimensional data
				\item Use cross-validation to select $\lambda$
				\item Start with no regularization, add if needed
				\item Balance bias and variance
			\end{itemize}
		\end{itemize}
		
		\vspace{0.05cm}
		\textcolor{myred}{\textbf{WHY OTHER OPTIONS ARE WRONG:}}
		\begin{itemize}
			\item \textcolor{myred}{A:} Underfitting would be worsened by regularization
			\item \textcolor{myred}{C:} Too few parameters = underfitting
			\item \textcolor{myred}{D:} Perfectly clean data still may benefit from regularization
		\end{itemize}
		
		\textcolor{myred}{\textbf{EXAM TRAP:}} Use regularization when overfitting (high variance)!
	\end{frame}
	
	\begin{frame}{Q71: L1 vs L2 Feature Selection - Tutorial}
		\fontsize{6.5}{7.5}\selectfont
		\textbf{Which regularization method is more likely to perform feature selection?}
		
		\vspace{0.05cm}
		\textcolor{myblue}{\textbf{ANSWER: B}} - L1 regularization (Lasso)
		
		\vspace{0.05cm}
		\textbf{DETAILED EXPLANATION:}
		\begin{itemize}
			\item \textbf{Why L1 performs feature selection:}
			\begin{itemize}
				\item L1 penalty = $\lambda \sum |w_i|$
				\item Non-differentiable at zero
				\item Creates sparse solutions
				\item Weights exactly zero $\Rightarrow$ features removed
			\end{itemize}
			\item \textbf{Why L2 doesn't:}
			\begin{itemize}
				\item L2 penalty = $\lambda \sum w_i^2$
				\item Differentiable everywhere
				\item Shrinks weights but never to zero
				\item All features retained
			\end{itemize}
			\item \textbf{Geometric intuition:}
			\begin{itemize}
				\item L1 constraint: diamond shape
				\item L2 constraint: circle shape
				\item Diamond has corners on axes $\Rightarrow$ sparse solutions
			\end{itemize}
			\item \textbf{Practical implications:}
			\begin{itemize}
				\item Lasso for feature selection
				\item Ridge for shrinkage
				\item Elastic Net for both
			\end{itemize}
		\end{itemize}
		
		\vspace{0.05cm}
		\textcolor{myred}{\textbf{WHY OTHER OPTIONS ARE WRONG:}}
		\begin{itemize}
			\item \textcolor{myred}{A:} L2 only shrinks weights
			\item \textcolor{myred}{C:} They are not equal
			\item \textcolor{myred}{D:} L1 does perform feature selection
		\end{itemize}
		
		\textcolor{myred}{\textbf{EXAM TRAP:}} L1 (Lasso) = feature selection; L2 (Ridge) = weight shrinkage!
	\end{frame}
	
	\begin{frame}{Q72: Weight Zero Interpretation - Tutorial}
		\fontsize{6.5}{7.5}\selectfont
		\textbf{In regularization, what does a weight of zero indicate?}
		
		\vspace{0.05cm}
		\textcolor{myblue}{\textbf{ANSWER: B}} - The feature is not used by the model (feature selection).
		
		\vspace{0.05cm}
		\textbf{DETAILED EXPLANATION:}
		\begin{itemize}
			\item \textbf{Weight interpretation:}
			\begin{itemize}
				\item Weight $w_j$ multiplies feature $x_j$
				\item $w_j = 0 \Rightarrow$ feature $x_j$ contributes nothing
				\item Feature is effectively removed
			\end{itemize}
			\item \textbf{Why this happens:}
			\begin{itemize}
				\item L1 regularization encourages sparsity
				\item Feature is not predictive enough
				\item Feature is redundant
				\item Regularization selects better features
			\end{itemize}
			\item \textbf{Benefit:}
			\begin{itemize}
				\item Simpler model
				\item Better interpretability
				\item Reduced overfitting
				\item Faster prediction (fewer features)
			\end{itemize}
			\item \textbf{Caveat:}
			\begin{itemize}
				\item Zero weights don't mean feature has no predictive power
				\item Feature may be correlated with others
				\item Feature may be less important than others
			\end{itemize}
		\end{itemize}
		
		\vspace{0.05cm}
		\textcolor{myred}{\textbf{WHY OTHER OPTIONS ARE WRONG:}}
		\begin{itemize}
			\item \textcolor{myred}{A:} Zero weight means no importance
			\item \textcolor{myred}{C:} Weight doesn't indicate variance
			\item \textcolor{myred}{D:} Weight doesn't indicate bias
		\end{itemize}
		
		\textcolor{myred}{\textbf{EXAM TRAP:}} Weight = 0 $\to$ feature not used (feature selection)!
	\end{frame}
	
	\begin{frame}{Q73: $\lambda$ Too Large - Tutorial}
		\fontsize{6.5}{7.5}\selectfont
		\textbf{If $\lambda$ is too large in regularization, what happens?}
		
		\vspace{0.05cm}
		\textcolor{myblue}{\textbf{ANSWER: B}} - The model underfits.
		
		\vspace{0.05cm}
		\textbf{DETAILED EXPLANATION:}
		\begin{itemize}
			\item \textbf{Effect of large $\lambda$:}
			\begin{itemize}
				\item Penalty dominates the loss function
				\item All weights shrink toward zero
				\item Model becomes too simple
				\item Cannot capture data patterns
			\end{itemize}
			\item \textbf{Symptoms of underfitting:}
			\begin{itemize}
				\item High training error
				\item High test error
				\item Model predictions are too simple
				\item Weights all close to zero
			\end{itemize}
			\item \textbf{Comparison:}
			\begin{itemize}
				\item $\lambda = 0$: No regularization (may overfit)
				\item $\lambda$ small: Some regularization (balance)
				\item $\lambda$ large: Too much regularization (underfit)
				\item $\lambda \to \infty$: All weights $\to 0$ (extreme underfit)
			\end{itemize}
			\item \textbf{Fixing:}
			\begin{itemize}
				\item Decrease $\lambda$
				\item Use cross-validation to find optimal $\lambda$
				\item Start with small $\lambda$ and increase as needed
			\end{itemize}
		\end{itemize}
		
		\vspace{0.05cm}
		\textcolor{myred}{\textbf{WHY OTHER OPTIONS ARE WRONG:}}
		\begin{itemize}
			\item \textcolor{myred}{A:} Too small $\lambda$ leads to overfitting
			\item \textcolor{myred}{C:} Optimal requires balanced $\lambda$
			\item \textcolor{myred}{D:} $\lambda$ definitely affects model
		\end{itemize}
		
		\textcolor{myred}{\textbf{EXAM TRAP:}} Too large $\lambda$ $\to$ underfitting!
	\end{frame}
	
	\begin{frame}{Q74: Regularization Forms - Tutorial}
		\fontsize{6.5}{7.5}\selectfont
		\textbf{Which of the following is a form of regularization?}
		
		\vspace{0.05cm}
		\textcolor{myblue}{\textbf{ANSWER: D}} - All of the above (L1, L2, Dropout)
		
		\vspace{0.05cm}
		\textbf{DETAILED EXPLANATION:}
		\begin{itemize}
			\item \textbf{L1 regularization (Lasso):}
			\begin{itemize}
				\item Adds $\lambda \sum |w_i|$ to loss
				\item Performs feature selection
				\item Sparse solutions
			\end{itemize}
			\item \textbf{L2 regularization (Ridge):}
			\begin{itemize}
				\item Adds $\lambda \sum w_i^2$ to loss
				\item Shrinks weights
				\item Dense solutions
			\end{itemize}
			\item \textbf{Dropout:}
			\begin{itemize}
				\item Randomly drops neurons
				\item Ensemble effect
				\item Prevents co-adaptation
			\end{itemize}
			\item \textbf{Other forms not listed:}
			\begin{itemize}
				\item Early stopping
				\item Batch normalization (has regularization effect)
				\item Data augmentation
				\item Weight decay (same as L2 in many frameworks)
			\end{itemize}
		\end{itemize}
		
		\vspace{0.05cm}
		\textcolor{myred}{\textbf{WHY OTHER OPTIONS ARE WRONG:}}
		\begin{itemize}
			\item \textcolor{myred}{A:} Correct but incomplete
			\item \textcolor{myred}{B:} Correct but incomplete
			\item \textcolor{myred}{C:} Correct but incomplete
		\end{itemize}
		
		\textcolor{myred}{\textbf{EXAM TRAP:}} Many techniques are forms of regularization!
	\end{frame}
	
	\begin{frame}{Q75: $\lambda$ Purpose - Tutorial}
		\fontsize{6.5}{7.5}\selectfont
		\textbf{What is the purpose of the regularization parameter $\lambda$?}
		
		\vspace{0.05cm}
		\textcolor{myblue}{\textbf{ANSWER: A}} - To control the tradeoff between fitting the training data and keeping the model simple.
		
		\vspace{0.05cm}
		\textbf{DETAILED EXPLANATION:}
		\begin{itemize}
			\item \textbf{Role of $\lambda$:}
			\begin{itemize}
				\item Balances two competing objectives
				\item Fit training data well (small $\lambda$)
				\item Keep model simple (large $\lambda$)
				\item Tradeoff controlled by $\lambda$
			\end{itemize}
			\item \textbf{Extreme cases:}
			\begin{itemize}
				\item $\lambda = 0$: Minimizes only training error
				\item $\lambda \to \infty$: Minimizes only penalty (weights $\to 0$)
			\end{itemize}
			\item \textbf{Choosing $\lambda$:}
			\begin{itemize}
				\item Must be selected carefully
				\item Use cross-validation
				\item Tradeoff between bias and variance
				\item Problem-specific
			\end{itemize}
			\item \textbf{Relationship to other hyperparameters:}
			\begin{itemize}
				\item Different from learning rate ($\alpha$)
				\item Different from number of epochs
				\item Different from batch size
			\end{itemize}
		\end{itemize}
		
		\vspace{0.05cm}
		\textcolor{myred}{\textbf{WHY OTHER OPTIONS ARE WRONG:}}
		\begin{itemize}
			\item \textcolor{myred}{B:} Learning rate is $\alpha$
			\item \textcolor{myred}{C:} Number of epochs is different
			\item \textcolor{myred}{D:} Batch size is different
		\end{itemize}
		
		\textcolor{myred}{\textbf{EXAM TRAP:}} $\lambda$ controls the bias-variance tradeoff in regularization!
	\end{frame}
	
	% ============================================================
	% SECTION 6: CROSS-VALIDATION - QUESTIONS 76-90
	% ============================================================
	
	\begin{frame}{Q76: Cross-Validation Definition - Tutorial}
		\fontsize{6.5}{7.5}\selectfont
		\textbf{What is cross-validation in supervised learning?}
		
		\vspace{0.05cm}
		\textcolor{myblue}{\textbf{ANSWER: B}} - Resampling technique used to evaluate model performance and select hyperparameters.
		
		\vspace{0.05cm}
		\textbf{DETAILED EXPLANATION:}
		\begin{itemize}
			\item \textbf{Definition:} Technique for assessing model generalization
			\begin{itemize}
				\item Partitions data into training and validation sets
				\item Multiple partitions for reliable estimate
				\item Uses all data for both training and validation
			\end{itemize}
			\item \textbf{Key purposes:}
			\begin{itemize}
				\item \textbf{Model evaluation:} Estimate test performance
				\item \textbf{Hyperparameter tuning:} Select best parameters
				\item \textbf{Model comparison:} Choose between models
			\end{itemize}
			\item \textbf{Why it's needed:}
			\begin{itemize}
				\item Single train/test split is unreliable
				\item Limited data: want to use all for training
				\item Avoid overfitting to a specific split
			\end{itemize}
			\item \textbf{Common methods:}
			\begin{itemize}
				\item k-fold cross-validation
				\item Leave-one-out cross-validation (LOOCV)
				\item Stratified cross-validation
			\end{itemize}
		\end{itemize}
		
		\vspace{0.05cm}
		\textcolor{myred}{\textbf{WHY OTHER OPTIONS ARE WRONG:}}
		\begin{itemize}
			\item \textcolor{myred}{A:} This is data augmentation
			\item \textcolor{myred}{C:} This is regularization
			\item \textcolor{myred}{D:} This is initialization
		\end{itemize}
		
		\textcolor{myred}{\textbf{EXAM TRAP:}} Cross-validation evaluates model performance and selects hyperparameters!
	\end{frame}
	
	\begin{frame}{Q77: K-Fold Cross-Validation - Tutorial}
		\fontsize{6.5}{7.5}\selectfont
		\textbf{How does k-fold cross-validation work?}
		
		\vspace{0.05cm}
		\textcolor{myblue}{\textbf{ANSWER: B}} - Split into k parts, train on k-1 parts and validate on remaining 1 part, repeat k times.
		
		\vspace{0.05cm}
		\textbf{DETAILED EXPLANATION:}
		\begin{itemize}
			\item \textbf{Step-by-step:}
			\begin{enumerate}
				\item Randomly split data into $k$ folds of equal size
				\item For $i = 1$ to $k$:
				\begin{itemize}
					\item Use fold $i$ as validation set
					\item Use remaining $k-1$ folds as training set
					\item Train model and evaluate
				\end{itemize}
				\item Average performance over all $k$ iterations
			\end{enumerate}
			\item \textbf{Example (k=5):}
			\begin{itemize}
				\item 1000 samples, 5 folds of 200 each
				\item Iteration 1: Train on 800, Validate on 200
				\item Repeat 5 times
				\item Average of 5 validation scores
			\end{itemize}
			\item \textbf{Advantages:}
			\begin{itemize}
				\item All data used for training (k-1 times)
				\item All data used for validation (once)
				\item Reliable performance estimate
				\item Reduces overfitting to validation set
			\end{itemize}
			\item \textbf{Common choices:}
			\begin{itemize}
				\item k=5, k=10 (most common)
				\item k=3 for very small datasets
				\item k=10 for moderate datasets
			\end{itemize}
		\end{itemize}
		
		\vspace{0.05cm}
		\textcolor{myred}{\textbf{WHY OTHER OPTIONS ARE WRONG:}}
		\begin{itemize}
			\item \textcolor{myred}{A:} Reverses training and validation
			\item \textcolor{myred}{C:} All training = no validation
			\item \textcolor{myred}{D:} All validation = no training
		\end{itemize}
		
		\textcolor{myred}{\textbf{EXAM TRAP:}} K-fold CV = train on k-1 folds, validate on 1 fold, repeat k times!
	\end{frame}
	
	\begin{frame}{Q78: Stratified Cross-Validation - Tutorial}
		\fontsize{6.5}{7.5}\selectfont
		\textbf{What is stratified cross-validation?}
		
		\vspace{0.05cm}
		\textcolor{myblue}{\textbf{ANSWER: B}} - Each fold preserves the same class distribution as the full dataset.
		
		\vspace{0.05cm}
		\textbf{DETAILED EXPLANATION:}
		\begin{itemize}
			\item \textbf{Definition:} CV with class balance preservation
			\begin{itemize}
				\item Each fold has same proportion of each class
				\item Matches the overall class distribution
				\item Ensures representative validation sets
			\end{itemize}
			\item \textbf{Example:}
			\begin{itemize}
				\item Dataset: 90\% class A, 10\% class B
				\item k=5 folds: each fold has 90\% A, 10\% B
				\item Without stratification: some folds might have no B
			\end{itemize}
			\item \textbf{When to use:}
			\begin{itemize}
				\item Imbalanced datasets
				\item Classification problems
				\item When class proportions matter
			\end{itemize}
			\item \textbf{Benefits:}
			\begin{itemize}
				\item More reliable performance estimates
				\item Avoids extreme validation sets
				\item Better for rare classes
				\item Reduces variance of estimates
			\end{itemize}
		\end{itemize}
		
		\vspace{0.05cm}
		\textcolor{myred}{\textbf{WHY OTHER OPTIONS ARE WRONG:}}
		\begin{itemize}
			\item \textcolor{myred}{A:} Equal fold sizes is standard k-fold
			\item \textcolor{myred}{C:} Random folds may not preserve distribution
			\item \textcolor{myred}{D:} Sequential folds may introduce bias
		\end{itemize}
		
		\textcolor{myred}{\textbf{EXAM TRAP:}} Stratified CV preserves class proportions in each fold!
	\end{frame}
	
	\begin{frame}{Q79: Stratified CV - When to Use - Tutorial}
		\fontsize{6.5}{7.5}\selectfont
		\textbf{When is stratified cross-validation particularly important?}
		
		\vspace{0.05cm}
		\textcolor{myblue}{\textbf{ANSWER: B}} - When the dataset is imbalanced (one class is much rarer than others).
		
		\vspace{0.05cm}
		\textbf{DETAILED EXPLANATION:}
		\begin{itemize}
			\item \textbf{Imbalanced dataset example:}
			\begin{itemize}
				\item 95\% negative class, 5\% positive class
				\item Without stratification: some folds may have 0\% positive
				\item Validation performance varies wildly
				\item Estimates are unreliable
			\end{itemize}
			\item \textbf{Why it's important:}
			\begin{itemize}
				\item Rare class performance matters (fraud detection)
				\item AUC, F1 need proper class representation
				\item Avoids optimistic performance estimates
				\item Better model selection
			\end{itemize}
			\item \textbf{Without stratification:}
			\begin{itemize}
				\item Fold may have no rare class samples
				\item Cannot evaluate rare class performance
				\item Model selection may choose wrong model
			\end{itemize}
			\item \textbf{Practical advice:}
			\begin{itemize}
				\item Always use stratification for classification
				\item Especially when class imbalance exists
				\item Also consider weighted metrics
			\end{itemize}
		\end{itemize}
		
		\vspace{0.05cm}
		\textcolor{myred}{\textbf{WHY OTHER OPTIONS ARE WRONG:}}
		\begin{itemize}
			\item \textcolor{myred}{A:} Balanced datasets don't need stratification
			\item \textcolor{myred}{C:} Size doesn't determine need
			\item \textcolor{myred}{D:} Stratification is about class distribution
		\end{itemize}
		
		\textcolor{myred}{\textbf{EXAM TRAP:}} Stratified CV is essential for imbalanced datasets!
	\end{frame}
	
	\begin{frame}{Q80: LOOCV Definition - Tutorial}
		\fontsize{6.5}{7.5}\selectfont
		\textbf{What is Leave-One-Out Cross-Validation (LOOCV)?}
		
		\vspace{0.05cm}
		\textcolor{myblue}{\textbf{ANSWER: A}} - Special case of k-fold CV where k = N (one observation per fold).
		
		\vspace{0.05cm}
		\textbf{DETAILED EXPLANATION:}
		\begin{itemize}
			\item \textbf{Definition:} Extreme case of k-fold CV
			\begin{itemize}
				\item $k = N$ where $N$ is number of observations
				\item Each fold contains exactly one observation
				\item Train on $N-1$ observations each time
			\end{itemize}
			\item \textbf{Algorithm:}
			\begin{enumerate}
				\item For $i = 1$ to $N$:
				\begin{itemize}
					\item Hold out observation $i$ as validation
					\item Train on remaining $N-1$ observations
					\item Evaluate on held-out observation
				\end{itemize}
				\item Average over all $N$ iterations
			\end{enumerate}
			\item \textbf{Example:}
			\begin{itemize}
				\item N = 100 observations
				\item Train 100 models, each on 99 observations
				\item Each validation set has 1 observation
				\item Average validation error
			\end{itemize}
			\item \textbf{Properties:}
			\begin{itemize}
				\item Uses almost all data for training
				\item Very low bias
				\item High variance (estimates are correlated)
				\item Computationally expensive
			\end{itemize}
		\end{itemize}
		
		\vspace{0.05cm}
		\textcolor{myred}{\textbf{WHY OTHER OPTIONS ARE WRONG:}}
		\begin{itemize}
			\item \textcolor{myred}{B:} k=2 is 2-fold CV
			\item \textcolor{myred}{C:} k=10 is 10-fold CV
			\item \textcolor{myred}{D:} LOOCV uses validation sets
		\end{itemize}
		
		\textcolor{myred}{\textbf{EXAM TRAP:}} LOOCV = k-fold CV with k = N!
	\end{frame}
	
	\begin{frame}{Q81: LOOCV Advantages - Tutorial}
		\fontsize{6.5}{7.5}\selectfont
		\textbf{What is an advantage of Leave-One-Out Cross-Validation?}
		
		\vspace{0.05cm}
		\textcolor{myblue}{\textbf{ANSWER: B}} - Uses almost all data for training each time (N-1 observations), giving low bias.
		
		\vspace{0.05cm}
		\textbf{DETAILED EXPLANATION:}
		\begin{itemize}
			\item \textbf{Low bias advantage:}
			\begin{itemize}
				\item Training set size is $N-1$ (almost full dataset)
				\item Model sees almost all data each time
				\item Training set is representative
				\item No data withheld from training
			\end{itemize}
			\item \textbf{Why bias is low:}
			\begin{itemize}
				\item Model trained on nearly full dataset
				\item No loss from using less data
				\item Consistent with full data model
				\item Less training-test discrepancy
			\end{itemize}
			\item \textbf{Comparison to other CV:}
			\begin{itemize}
				\item 10-fold CV: trains on 90\% of data
				\item 5-fold CV: trains on 80\% of data
				\item LOOCV: trains on 99\%+ of data
				\item LOOCV has lowest bias
			\end{itemize}
			\item \textbf{When this matters:}
			\begin{itemize}
				\item Small datasets (every observation counts)
				\item When bias is a concern
				\item When you want best estimate
			\end{itemize}
		\end{itemize}
		
		\vspace{0.05cm}
		\textcolor{myred}{\textbf{WHY OTHER OPTIONS ARE WRONG:}}
		\begin{itemize}
			\item \textcolor{myred}{A:} LOOCV is computationally expensive
			\item \textcolor{myred}{C:} LOOCV can have high variance
			\item \textcolor{myred}{D:} LOOCV is not suitable for large datasets
		\end{itemize}
		
		\textcolor{myred}{\textbf{EXAM TRAP:}} LOOCV has low bias but is computationally expensive!
	\end{frame}
	
	\begin{frame}{Q82: LOOCV Disadvantages - Tutorial}
		\fontsize{6.5}{7.5}\selectfont
		\textbf{What is a disadvantage of Leave-One-Out Cross-Validation?}
		
		\vspace{0.05cm}
		\textcolor{myblue}{\textbf{ANSWER: B}} - It is computationally expensive (requires N training runs).
		
		\vspace{0.05cm}
		\textbf{DETAILED EXPLANATION:}
		\begin{itemize}
			\item \textbf{Computational cost:}
			\begin{itemize}
				\item Need to train model $N$ times
				\item Each training is on $N-1$ observations
				\item Total operations: $N \times \text{cost\_per\_train}$
				\item For $N = 10,000$: 10,000 training runs
			\end{itemize}
			\item \textbf{When it's problematic:}
			\begin{itemize}
				\item Large datasets ($N > 1000$)
				\item Complex models (slow to train)
				\item Limited computational resources
			\end{itemize}
			\item \textbf{Alternative solutions:}
			\begin{itemize}
				\item Use k-fold with smaller $k$ (e.g., 5 or 10)
				\item Use approximate LOOCV formulas (for some models)
				\item Use efficient implementations when available
			\end{itemize}
			\item \textbf{Other disadvantages:}
			\begin{itemize}
				\item High variance (estimates are correlated)
				\item Can be unstable
				\item Not suitable for large datasets
			\end{itemize}
		\end{itemize}
		
		\vspace{0.05cm}
		\textcolor{myred}{\textbf{WHY OTHER OPTIONS ARE WRONG:}}
		\begin{itemize}
			\item \textcolor{myred}{A:} LOOCV has low bias, not high bias
			\item \textcolor{myred}{C:} LOOCV doesn't necessarily overfit
			\item \textcolor{myred}{D:} LOOCV can be used for classification
		\end{itemize}
		
		\textcolor{myred}{\textbf{EXAM TRAP:}} LOOCV is computationally expensive for large N!
	\end{frame}
	
	\begin{frame}{Q83: CV Bias-Variance Tradeoff - Tutorial}
		\fontsize{6.5}{7.5}\selectfont
		\textbf{How does choice of k in k-fold CV affect bias-variance tradeoff?}
		
		\vspace{0.05cm}
		\textcolor{myblue}{\textbf{ANSWER: B}} - Larger k decreases bias and increases variance.
		
		\vspace{0.05cm}
		\textbf{DETAILED EXPLANATION:}
		\begin{itemize}
			\item \textbf{Effect of larger k:}
			\begin{itemize}
				\item More data used for training ($\frac{k-1}{k}$ fraction)
				\item Training sets more similar to full dataset
				\item Lower bias (models are better)
				\item Higher variance (estimates are more correlated)
			\end{itemize}
			\item \textbf{Effect of smaller k:}
			\begin{itemize}
				\item Less data used for training
				\item Training sets less representative
				\item Higher bias (models are worse)
				\item Lower variance (estimates are less correlated)
			\end{itemize}
			\item \textbf{The tradeoff:}
			\begin{itemize}
				\item LOOCV (k=N): Lowest bias, highest variance
				\item k=2: Highest bias, lowest variance
				\item k=10: Good balance (most common)
			\end{itemize}
			\item \textbf{Choosing k:}
			\begin{itemize}
				\item k=10 is standard
				\item k=5 for larger datasets
				\item k=N (LOOCV) for small datasets
			\end{itemize}
		\end{itemize}
		
		\vspace{0.05cm}
		\textcolor{myred}{\textbf{WHY OTHER OPTIONS ARE WRONG:}}
		\begin{itemize}
			\item \textcolor{myred}{A:} This is the opposite direction
			\item \textcolor{myred}{C:} k definitely affects bias-variance
			\item \textcolor{myred}{D:} Smaller k increases bias, decreases variance
		\end{itemize}
		
		\textcolor{myred}{\textbf{EXAM TRAP:}} Larger k $\to$ lower bias, higher variance!
	\end{frame}
	
	\begin{frame}{Q84: Common k Value - Tutorial}
		\fontsize{6.5}{7.5}\selectfont
		\textbf{What is the most common value of k in k-fold cross-validation?}
		
		\vspace{0.05cm}
		\textcolor{myblue}{\textbf{ANSWER: C}} - k = 10
		
		\vspace{0.05cm}
		\textbf{DETAILED EXPLANATION:}
		\begin{itemize}
			\item \textbf{Why k=10 is common:}
			\begin{itemize}
				\item Good balance between bias and variance
				\item Computationally feasible
				\item Widely studied and recommended
				\item Standard in practice
			\end{itemize}
			\item \textbf{Other common values:}
			\begin{itemize}
				\item \textbf{k=5:} For larger datasets
				\item \textbf{k=10:} Standard default
				\item \textbf{k=LOOCV:} For small datasets
			\end{itemize}
			\item \textbf{Rules of thumb:}
			\begin{itemize}
				\item Use k=10 for most problems
				\item Use k=5 for $N > 10,000$
				\item Use LOOCV for $N < 100$
				\item Increase k for smaller datasets
			\end{itemize}
			\item \textbf{Tradeoffs for different k:}
			\begin{itemize}
				\item k=10: $90\%$ training, $10\%$ validation
				\item k=5: $80\%$ training, $20\%$ validation
				\item k=2: $50\%$ training, $50\%$ validation (too little training)
			\end{itemize}
		\end{itemize}
		
		\vspace{0.05cm}
		\textcolor{myred}{\textbf{WHY OTHER OPTIONS ARE WRONG:}}
		\begin{itemize}
			\item \textcolor{myred}{A:} k=2 is sometimes used but not most common
			\item \textcolor{myred}{B:} k=5 is also common but k=10 is most common
			\item \textcolor{myred}{D:} LOOCV is used for small datasets
		\end{itemize}
		
		\textcolor{myred}{\textbf{EXAM TRAP:}} k = 10 is the most common value!
	\end{frame}
	
	\begin{frame}{Q85: CV Purpose - Tutorial}
		\fontsize{6.5}{7.5}\selectfont
		\textbf{What is the primary purpose of cross-validation?}
		
		\vspace{0.05cm}
		\textcolor{myblue}{\textbf{ANSWER: B}} - To evaluate how well a model will generalize to new, unseen data.
		
		\vspace{0.05cm}
		\textbf{DETAILED EXPLANATION:}
		\begin{itemize}
			\item \textbf{Primary goal:} Estimate generalization performance
			\begin{itemize}
				\item How well does model perform on new data?
				\item What is expected test error?
				\item Can the model be trusted?
			\end{itemize}
			\item \textbf{Why we need it:}
			\begin{itemize}
				\item Training performance is optimistic
				\item Test data is not available during development
				\item Need reliable estimate without overfitting
			\end{itemize}
			\item \textbf{Secondary purposes:}
			\begin{itemize}
				\item Hyperparameter selection
				\item Model comparison
				\item Feature selection
				\item Ensemble creation
			\end{itemize}
			\item \textbf{What it doesn't do:}
			\begin{itemize}
				\item Increase training data (that's data augmentation)
				\item Reduce overfitting (that's regularization)
				\item Select features (though it helps)
			\end{itemize}
		\end{itemize}
		
		\vspace{0.05cm}
		\textcolor{myred}{\textbf{WHY OTHER OPTIONS ARE WRONG:}}
		\begin{itemize}
			\item \textcolor{myred}{A:} This is data augmentation
			\item \textcolor{myred}{C:} This is feature selection
			\item \textcolor{myred}{D:} This is a training parameter
		\end{itemize}
		
		\textcolor{myred}{\textbf{EXAM TRAP:}} Cross-validation estimates generalization performance!
	\end{frame}
	
	\begin{frame}{Q86: Bootstrapping Definition - Tutorial}
		\fontsize{6.5}{7.5}\selectfont
		\textbf{What is bootstrapping in model evaluation?}
		
		\vspace{0.05cm}
		\textcolor{myblue}{\textbf{ANSWER: A}} - Resampling technique that samples with replacement from the dataset.
		
		\vspace{0.05cm}
		\textbf{DETAILED EXPLANATION:}
		\begin{itemize}
			\item \textbf{Definition:} Sampling with replacement
			\begin{itemize}
				\item Create $B$ bootstrap samples
				\item Each sample size = $N$ (same as original)
				\item Sampling with replacement
				\item Some observations appear multiple times
				\item Some observations don't appear
			\end{itemize}
			\item \textbf{Expected composition:}
			\begin{itemize}
				\item Each bootstrap sample has $N$ observations
				\item Approximately $63.2\%$ of original data appears
				\item Remaining $36.8\%$ are out-of-bag (OOB)
			\end{itemize}
			\item \textbf{Usage:}
			\begin{itemize}
				\item Estimate variability of model performance
				\item Construct confidence intervals
				\item Evaluate stability
				\item Used in bagging (Bootstrap Aggregating)
			\end{itemize}
			\item \textbf{Example:}
			\begin{itemize}
				\item Original: [1, 2, 3, 4, 5]
				\item Bootstrap: [1, 3, 3, 4, 5] (3 appears twice)
				\item Another: [1, 2, 2, 4, 4] (2 and 4 appear twice)
			\end{itemize}
		\end{itemize}
		
		\vspace{0.05cm}
		\textcolor{myred}{\textbf{WHY OTHER OPTIONS ARE WRONG:}}
		\begin{itemize}
			\item \textcolor{myred}{B:} Sampling without replacement is subsampling
			\item \textcolor{myred}{C:} Bootstrapping is not regularization
			\item \textcolor{myred}{D:} Bootstrapping is not feature selection
		\end{itemize}
		
		\textcolor{myred}{\textbf{EXAM TRAP:}} Bootstrapping = sampling with replacement!
	\end{frame}
	
	\begin{frame}{Q87: Bootstrapping Purpose - Tutorial}
		\fontsize{6.5}{7.5}\selectfont
		\textbf{What is the primary purpose of bootstrapping in model evaluation?}
		
		\vspace{0.05cm}
		\textcolor{myblue}{\textbf{ANSWER: A}} - To estimate the variability of model performance (e.g., confidence intervals).
		
		\vspace{0.05cm}
		\textbf{DETAILED EXPLANATION:}
		\begin{itemize}
			\item \textbf{Primary goal:} Quantify uncertainty
			\begin{itemize}
				\item Model performance varies across samples
				\item Need to know how much it varies
				\item Construct confidence intervals
				\item Assess model stability
			\end{itemize}
			\item \textbf{How it works:}
			\begin{itemize}
				\item Create $B$ bootstrap samples
				\item Train model on each sample
				\item Evaluate performance on each sample (or OOB)
				\item Get distribution of performance metrics
			\end{itemize}
			\item \textbf{Applications:}
			\begin{itemize}
				\item Confidence intervals for accuracy
				\item Confidence intervals for coefficients
				\item Model stability assessment
				\item Parameter uncertainty quantification
			\end{itemize}
			\item \textbf{Comparison to CV:}
			\begin{itemize}
				\item CV: estimates average performance
				\item Bootstrap: estimates performance distribution
				\item Both useful for different purposes
			\end{itemize}
		\end{itemize}
		
		\vspace{0.05cm}
		\textcolor{myred}{\textbf{WHY OTHER OPTIONS ARE WRONG:}}
		\begin{itemize}
			\item \textcolor{myred}{B:} This is regularization
			\item \textcolor{myred}{C:} Learning rate is a training parameter
			\item \textcolor{myred}{D:} This is feature selection
		\end{itemize}
		
		\textcolor{myred}{\textbf{EXAM TRAP:}} Bootstrapping estimates performance variability!
	\end{frame}
	
	\begin{frame}{Q88: CV vs Bootstrapping - Tutorial}
		\fontsize{6.5}{7.5}\selectfont
		\textbf{What is the key difference between cross-validation and bootstrapping?}
		
		\vspace{0.05cm}
		\textcolor{myblue}{\textbf{ANSWER: A}} - Cross-validation samples without replacement; bootstrapping samples with replacement.
		
		\vspace{0.05cm}
		\textbf{DETAILED EXPLANATION:}
		\begin{itemize}
			\item \textbf{Cross-validation:}
			\begin{itemize}
				\item Samples without replacement
				\item Each observation appears exactly once in validation
				\item Partitions are disjoint
				\item Training sets are subsets of data
				\item No overlap between folds
			\end{itemize}
			\item \textbf{Bootstrapping:}
			\begin{itemize}
				\item Samples with replacement
				\item Observations can appear multiple times
				\item Each sample is same size as original
				\item Some observations omitted (OOB)
				\item Overlap between samples
			\end{itemize}
			\item \textbf{Other differences:}
			\begin{itemize}
				\item CV: deterministic (for given random seed)
				\item Bootstrap: random resampling
				\item CV: evaluates on held-out data
				\item Bootstrap: evaluates on sample or OOB
			\end{itemize}
			\item \textbf{When to use each:}
			\begin{itemize}
				\item CV: model selection, hyperparameter tuning
				\item Bootstrap: uncertainty quantification, ensemble methods
			\end{itemize}
		\end{itemize}
		
		\vspace{0.05cm}
		\textcolor{myred}{\textbf{WHY OTHER OPTIONS ARE WRONG:}}
		\begin{itemize}
			\item \textcolor{myred}{B:} Both use parts of data
			\item \textcolor{myred}{C:} Both can be used for either task
			\item \textcolor{myred}{D:} They are different techniques
		\end{itemize}
		
		\textcolor{myred}{\textbf{EXAM TRAP:}} CV = without replacement; Bootstrapping = with replacement!
	\end{frame}
	
	\begin{frame}{Q89: Grid Search Definition - Tutorial}
		\fontsize{6.5}{7.5}\selectfont
		\textbf{What is grid search in hyperparameter tuning?}
		
		\vspace{0.05cm}
		\textcolor{myblue}{\textbf{ANSWER: B}} - Method that exhaustively searches over a predefined grid of hyperparameter values.
		
		\vspace{0.05cm}
		\textbf{DETAILED EXPLANATION:}
		\begin{itemize}
			\item \textbf{Definition:} Exhaustive parameter search
			\begin{itemize}
				\item Define grid of values for each hyperparameter
				\item Evaluate all combinations
				\item Select best combination
				\item Often combined with cross-validation
			\end{itemize}
			\item \textbf{Example:}
			\begin{itemize}
				\item $\lambda$: [0.001, 0.01, 0.1, 1]
				\item Learning rate: [0.001, 0.01, 0.1]
				\item Total combinations: $4 \times 3 = 12$
				\item Each combination trained and evaluated
			\end{itemize}
			\item \textbf{Advantages:}
			\begin{itemize}
				\item Guarantees finding best in grid
				\item Simple and easy to implement
				\item Reproducible results
				\item Parallelizable
			\end{itemize}
			\item \textbf{Disadvantages:}
			\begin{itemize}
				\item Computationally expensive
				\item Curse of dimensionality
				\item Continuous parameters difficult
				\item May miss good values
			\end{itemize}
		\end{itemize}
		
		\vspace{0.05cm}
		\textcolor{myred}{\textbf{WHY OTHER OPTIONS ARE WRONG:}}
		\begin{itemize}
			\item \textcolor{myred}{A:} This is random search
			\item \textcolor{myred}{C:} This is gradient-based optimization
			\item \textcolor{myred}{D:} Cross-validation is used WITH grid search
		\end{itemize}
		
		\textcolor{myred}{\textbf{EXAM TRAP:}} Grid search = exhaustive search over parameter grid!
	\end{frame}
	
	\begin{frame}{Q90: Grid Search Cost - Tutorial}
		\fontsize{6.5}{7.5}\selectfont
		\textbf{What is a major disadvantage of grid search?}
		
		\vspace{0.05cm}
		\textcolor{myblue}{\textbf{ANSWER: A}} - It is computationally expensive for large grids.
		
		\vspace{0.05cm}
		\textbf{DETAILED EXPLANATION:}
		\begin{itemize}
			\item \textbf{Computational complexity:}
			\begin{itemize}
				\item $M$ hyperparameters, each with $N_i$ values
				\item Total combinations: $\prod_{i=1}^M N_i$
				\item Grows exponentially with number of hyperparameters
				\item $M=5$, $N=10$ each $\Rightarrow$ 100,000 combinations
			\end{itemize}
			\item \textbf{When it's problematic:}
			\begin{itemize}
				\item Many hyperparameters
				\item Large grids for each parameter
				\item Expensive models to train
				\item Combined with cross-validation (multiplies cost)
			\end{itemize}
			\item \textbf{Alternatives:}
			\begin{itemize}
				\item Random search: sample random combinations
				\item Bayesian optimization: informed search
				\item Automated hyperparameter tuning
				\item Smaller grids with coarse-to-fine search
			\end{itemize}
			\item \textbf{Practical tips:}
			\begin{itemize}
				\item Start with coarse grid
				\item Use random search for many parameters
				\item Use Bayesian optimization for expensive models
			\end{itemize}
		\end{itemize}
		
		\vspace{0.05cm}
		\textcolor{myred}{\textbf{WHY OTHER OPTIONS ARE WRONG:}}
		\begin{itemize}
			\item \textcolor{myred}{B:} Finds best within grid (not suboptimal)
			\item \textcolor{myred}{C:} Can be used with CV
			\item \textcolor{myred}{D:} Doesn't require gradients
		\end{itemize}
		
		\textcolor{myred}{\textbf{EXAM TRAP:}} Grid search is computationally expensive for large grids!
	\end{frame}
	
	% ============================================================
	% SECTION 7: ACCURACY-INTERPRETABILITY - QUESTIONS 91-100
	% ============================================================
	
	\begin{frame}{Q91: Accuracy-Interpretability Tradeoff - Tutorial}
		\fontsize{6.5}{7.5}\selectfont
		\textbf{What is the accuracy-interpretability tradeoff?}
		
		\vspace{0.05cm}
		\textcolor{myblue}{\textbf{ANSWER: B}} - Tradeoff between predictive performance and how easily decisions can be understood.
		
		\vspace{0.05cm}
		\textbf{DETAILED EXPLANATION:}
		\begin{itemize}
			\item \textbf{Definition:} Tension between two competing goals
			\begin{itemize}
				\item \textbf{Accuracy:} How well the model predicts
				\item \textbf{Interpretability:} How well humans understand the model
				\item Often in conflict
			\end{itemize}
			\item \textbf{Why they conflict:}
			\begin{itemize}
				\item Simpler models are more interpretable
				\item Complex models are more accurate
				\item Tradeoff is fundamental
				\item Cannot maximize both
			\end{itemize}
			\item \textbf{Examples:}
			\begin{itemize}
				\item Linear regression: Very interpretable, limited accuracy
				\item Random forest: Moderate interpretability, high accuracy
				\item Deep neural network: Low interpretability, very high accuracy
			\end{itemize}
			\item \textbf{Choosing the tradeoff:}
			\begin{itemize}
				\item Depends on application
				\item Regulated industries: prefer interpretability
				\item Performance-critical: prefer accuracy
			\end{itemize}
		\end{itemize}
		
		\vspace{0.05cm}
		\textcolor{myred}{\textbf{WHY OTHER OPTIONS ARE WRONG:}}
		\begin{itemize}
			\item \textcolor{myred}{A:} This is speed-complexity tradeoff
			\item \textcolor{myred}{C:} This is bias-variance tradeoff
			\item \textcolor{myred}{D:} This is curse of dimensionality
		\end{itemize}
		
		\textcolor{myred}{\textbf{EXAM TRAP:}} Accuracy vs Interpretability = performance vs understandability!
	\end{frame}
	
	\begin{frame}{Q92: High Interpretability Models - Tutorial}
		\fontsize{6.5}{7.5}\selectfont
		\textbf{Which model typically has high interpretability?}
		
		\vspace{0.05cm}
		\textcolor{myblue}{\textbf{ANSWER: C}} - Decision trees (shallow)
		
		\vspace{0.05cm}
		\textbf{DETAILED EXPLANATION:}
		\begin{itemize}
			\item \textbf{Why decision trees are interpretable:}
			\begin{itemize}
				\item Decision rules are easy to follow
				\item Can be visualized
				\item Each split corresponds to a rule
				\item Path from root to leaf explains prediction
			\end{itemize}
			\item \textbf{Example:}
			\begin{itemize}
				\item If age > 30: predict high risk
				\item If age $\leq$ 30 and income > 50k: predict low risk
				\item Each prediction can be explained
			\end{itemize}
			\item \textbf{Other interpretable models:}
			\begin{itemize}
				\item Linear regression (coefficients explain)
				\item Logistic regression (odds ratios)
				\item Decision rules
				\item Naive Bayes
			\end{itemize}
			\item \textbf{Why others are less interpretable:}
			\begin{itemize}
				\item DNNs: black-box, millions of parameters
				\item Random forests: ensemble of many trees
				\item GBM: many trees, complex interactions
			\end{itemize}
		\end{itemize}
		
		\vspace{0.05cm}
		\textcolor{myred}{\textbf{WHY OTHER OPTIONS ARE WRONG:}}
		\begin{itemize}
			\item \textcolor{myred}{A:} DNNs are black-box models
			\item \textcolor{myred}{B:} Random forests are less interpretable
			\item \textcolor{myred}{D:} GBM is complex ensemble
		\end{itemize}
		
		\textcolor{myred}{\textbf{EXAM TRAP:}} Decision trees are highly interpretable "white-box" models!
	\end{frame}
	
	\begin{frame}{Q93: High Accuracy Models - Tutorial}
		\fontsize{6.5}{7.5}\selectfont
		\textbf{Which model typically has high accuracy but low interpretability?}
		
		\vspace{0.05cm}
		\textcolor{myblue}{\textbf{ANSWER: C}} - Deep neural networks
		
		\vspace{0.05cm}
		\textbf{DETAILED EXPLANATION:}
		\begin{itemize}
			\item \textbf{Why DNNs are accurate:}
			\begin{itemize}
				\item Can approximate any function
				\item Learn hierarchical features
				\item Massive parameter count
				\item State-of-the-art on many tasks
			\end{itemize}
			\item \textbf{Why DNNs are not interpretable:}
			\begin{itemize}
				\item Millions of weights
				\item No simple decision rules
				\item Features are abstract
				\item "Black-box" nature
				\item Hard to explain individual predictions
			\end{itemize}
			\item \textbf{Examples of high accuracy tasks:}
			\begin{itemize}
				\item Image recognition (95\%+ accuracy)
				\item Speech recognition
				\item Natural language processing
				\item Game playing (AlphaGo)
			\end{itemize}
			\item \textbf{Efforts to interpret DNNs:}
			\begin{itemize}
				\item SHAP values
				\item LIME
				\item Saliency maps
				\item Feature visualization
			\end{itemize}
		\end{itemize}
		
		\vspace{0.05cm}
		\textcolor{myred}{\textbf{WHY OTHER OPTIONS ARE WRONG:}}
		\begin{itemize}
			\item \textcolor{myred}{A:} Linear regression is highly interpretable
			\item \textcolor{myred}{B:} Decision trees are highly interpretable
			\item \textcolor{myred}{D:} Logistic regression is highly interpretable
		\end{itemize}
		
		\textcolor{myred}{\textbf{EXAM TRAP:}} Deep neural networks = high accuracy, low interpretability!
	\end{frame}
	
	\begin{frame}{Q94: Interpretability Importance - Tutorial}
		\fontsize{6.5}{7.5}\selectfont
		\textbf{In which domain is interpretability particularly important?}
		
		\vspace{0.05cm}
		\textcolor{myblue}{\textbf{ANSWER: C}} - Risk management and regulated industries (e.g., credit scoring)
		
		\vspace{0.05cm}
		\textbf{DETAILED EXPLANATION:}
		\begin{itemize}
			\item \textbf{Why regulated industries need interpretability:}
			\begin{itemize}
				\item Legal requirements: must explain decisions
				\item Regulatory compliance (GDPR, etc.)
				\item Right to explanation for consumers
				\item Avoiding discrimination
				\item Auditing and accountability
			\end{itemize}
			\item \textbf{Specific examples:}
			\begin{itemize}
				\item \textbf{Credit scoring:} Must explain why loan denied
				\item \textbf{Healthcare:} Must justify treatment decisions
				\item \textbf{Insurance:} Must explain premium calculations
				\item \textbf{Fraud detection:} Must justify alerts
			\end{itemize}
			\item \textbf{Why other domains differ:}
			\begin{itemize}
				\item Image recognition: interpretability less critical
				\item NLP: accuracy often prioritized
				\item Recommendations: personalization prioritized
			\end{itemize}
			\item \textbf{Balance in risk management:}
			\begin{itemize}
				\item White-box models preferred
				\item Even if accuracy is lower
				\item Explainability is legally required
			\end{itemize}
		\end{itemize}
		
		\vspace{0.05cm}
		\textcolor{myred}{\textbf{WHY OTHER OPTIONS ARE WRONG:}}
		\begin{itemize}
			\item \textcolor{myred}{A:} Image recognition prioritizes accuracy
			\item \textcolor{myred}{B:} NLP often prioritizes accuracy
			\item \textcolor{myred}{D:} Recommendations prioritize personalization
		\end{itemize}
		
		\textcolor{myred}{\textbf{EXAM TRAP:}} Interpretability is critical in risk management and regulated industries!
	\end{frame}
	
	\begin{frame}{Q95: Risk Management Model Choice - Tutorial}
		\fontsize{6.5}{7.5}\selectfont
		\textbf{In risk management, which type of model is often preferred?}
		
		\vspace{0.05cm}
		\textcolor{myblue}{\textbf{ANSWER: B}} - White-box models (e.g., decision trees, logistic regression)
		
		\vspace{0.05cm}
		\textbf{DETAILED EXPLANATION:}
		\begin{itemize}
			\item \textbf{Why white-box models are preferred:}
			\begin{itemize}
				\item Regulatory requirements demand transparency
				\item Must explain decisions to customers
				\item Auditing is easier
				\item Legal defensibility
				\item Fairness can be verified
			\end{itemize}
			\item \textbf{White-box examples in risk management:}
			\begin{itemize}
				\item Logistic regression for credit scoring
				\item Decision trees for fraud detection
				\item Linear models for risk assessment
				\item Scorecards in banking
			\end{itemize}
			\item \textbf{Tradeoff accepted:}
			\begin{itemize}
				\item May sacrifice some accuracy
				\item Interpretability is worth the cost
				\item Regulatory compliance is non-negotiable
			\end{itemize}
			\item \textbf{When black-box models are acceptable:}
			\begin{itemize}
				\item Internal research
				\item Non-regulated applications
				\item When explainability is not required
			\end{itemize}
		\end{itemize}
		
		\vspace{0.05cm}
		\textcolor{myred}{\textbf{WHY OTHER OPTIONS ARE WRONG:}}
		\begin{itemize}
			\item \textcolor{myred}{A:} Black-box models are avoided in risk management
			\item \textcolor{myred}{C:} Ensembles are less interpretable
			\item \textcolor{myred}{D:} Highest accuracy ignores interpretability requirements
		\end{itemize}
		
		\textcolor{myred}{\textbf{EXAM TRAP:}} Risk management prefers interpretable (white-box) models!
	\end{frame}
	
	\begin{frame}{Q96: Tradeoff Summary - Tutorial}
		\fontsize{6.5}{7.5}\selectfont
		\textbf{Which statement about the accuracy-interpretability tradeoff is TRUE?}
		
		\vspace{0.05cm}
		\textcolor{myblue}{\textbf{ANSWER: C}} - There is often a tradeoff between accuracy and interpretability.
		
		\vspace{0.05cm}
		\textbf{DETAILED EXPLANATION:}
		\begin{itemize}
			\item \textbf{Why there is a tradeoff:}
			\begin{itemize}
				\item Interpretability requires simplicity
				\item Accuracy often requires complexity
				\item These goals are in tension
				\item Cannot maximize both simultaneously
			\end{itemize}
			\item \textbf{Examples of the tradeoff:}
			\begin{itemize}
				\item Linear regression: Interpretable but limited
				\item Random forest: Less interpretable but more accurate
				\item DNN: Hard to interpret but highly accurate
			\end{itemize}
			\item \textbf{When the tradeoff is less severe:}
			\begin{itemize}
				\item Simple problems: both can be achieved
				\item Domain knowledge helps interpretability
				\item Some techniques improve both
			\end{itemize}
			\item \textbf{Managing the tradeoff:}
			\begin{itemize}
				\item Choose based on application
				\item Use interpretability tools (SHAP, LIME)
				\item Consider hybrid approaches
			\end{itemize}
		\end{itemize}
		
		\vspace{0.05cm}
		\textcolor{myred}{\textbf{WHY OTHER OPTIONS ARE WRONG:}}
		\begin{itemize}
			\item \textcolor{myred}{A:} False - simpler models often have lower accuracy
			\item \textcolor{myred}{B:} False - complex models often have higher accuracy
			\item \textcolor{myred}{D:} False - they are related through complexity
		\end{itemize}
		
		\textcolor{myred}{\textbf{EXAM TRAP:}} Accuracy and interpretability are often in tension (tradeoff)!
	\end{frame}
	
	\begin{frame}{Q97: Feature Importance - Tutorial}
		\fontsize{6.5}{7.5}\selectfont
		\textbf{What is feature importance in model interpretability?}
		
		\vspace{0.05cm}
		\textcolor{myblue}{\textbf{ANSWER: A}} - Measure of how important each feature is in making predictions.
		
		\vspace{0.05cm}
		\textbf{DETAILED EXPLANATION:}
		\begin{itemize}
			\item \textbf{Definition:} Quantifies feature contribution
			\begin{itemize}
				\item Which features drive predictions?
				\item How much does each feature contribute?
				\item Helps understand model behavior
				\item Identifies most important predictors
			\end{itemize}
			\item \textbf{Common methods:}
			\begin{itemize}
				\item Coefficient magnitude (linear models)
				\item Feature importance (tree-based models)
				\item Permutation importance
				\item SHAP values
				\item LIME weights
			\end{itemize}
			\item \textbf{Interpreting feature importance:}
			\begin{itemize}
				\item Higher value = more important
				\item Zero = no importance
				\item Positive vs negative direction matters
				\item Context-dependent
			\end{itemize}
			\item \textbf{Applications:}
			\begin{itemize}
				\item Feature selection
				\item Model explanation
				\item Domain insight
				\item Debugging models
			\end{itemize}
		\end{itemize}
		
		\vspace{0.05cm}
		\textcolor{myred}{\textbf{WHY OTHER OPTIONS ARE WRONG:}}
		\begin{itemize}
			\item \textcolor{myred}{B:} This is data availability
			\item \textcolor{myred}{C:} This is feature variance
			\item \textcolor{myred}{D:} This is correlation
		\end{itemize}
		
		\textcolor{myred}{\textbf{EXAM TRAP:}} Feature importance helps interpret model predictions!
	\end{frame}
	
	\begin{frame}{Q98: SHAP Values - Tutorial}
		\fontsize{6.5}{7.5}\selectfont
		\textbf{What are SHAP (SHapley Additive exPlanations) values?}
		
		\vspace{0.05cm}
		\textcolor{myblue}{\textbf{ANSWER: A}} - Method to compute feature importance using game theory.
		
		\vspace{0.05cm}
		\textbf{DETAILED EXPLANATION:}
		\begin{itemize}
			\item \textbf{Definition:} SHapley Additive exPlanations
			\begin{itemize}
				\item Based on Shapley values from game theory
				\item Distributes prediction among features
				\item Each feature gets a contribution
				\item Unifies other explanation methods
			\end{itemize}
			\item \textbf{Key properties:}
			\begin{itemize}
				\item \textbf{Additivity:} Sum of SHAP values = prediction
				\item \textbf{Consistency:} If model changes, values change accordingly
				\item \textbf{Local accuracy:} Explains individual predictions
				\item \textbf{Missingness:} Missing features get SHAP=0
			\end{itemize}
			\item \textbf{Advantages:}
			\begin{itemize}
				\item Globally and locally consistent
				\item Works for any model
				\item Provides both direction and magnitude
				\item Interpretable and reliable
			\end{itemize}
			\item \textbf{Example interpretation:}
			\begin{itemize}
				\item Base prediction: 0.5
				\item Feature A contribution: +0.3
				\item Feature B contribution: -0.1
				\item Final prediction: 0.7
			\end{itemize}
		\end{itemize}
		
		\vspace{0.05cm}
		\textcolor{myred}{\textbf{WHY OTHER OPTIONS ARE WRONG:}}
		\begin{itemize}
			\item \textcolor{myred}{B:} SHAP is for explanation, not accuracy
			\item \textcolor{myred}{C:} SHAP is not regularization
			\item \textcolor{myred}{D:} SHAP is not hyperparameter tuning
		\end{itemize}
		
		\textcolor{myred}{\textbf{EXAM TRAP:}} SHAP uses game theory to explain predictions!
	\end{frame}
	
	\begin{frame}{Q99: LIME - Tutorial}
		\fontsize{6.5}{7.5}\selectfont
		\textbf{What is LIME (Local Interpretable Model-agnostic Explanations)?}
		
		\vspace{0.05cm}
		\textcolor{myblue}{\textbf{ANSWER: A}} - Method to explain individual predictions by approximating model locally with interpretable model.
		
		\vspace{0.05cm}
		\textbf{DETAILED EXPLANATION:}
		\begin{itemize}
			\item \textbf{Definition:} Local Interpretable Model-agnostic Explanations
			\begin{itemize}
				\item Explains individual predictions
				\item Approximates complex model locally
				\item Uses interpretable model (linear, decision tree)
				\item Model-agnostic (works for any model)
			\end{itemize}
			\item \textbf{How it works:}
			\begin{enumerate}
				\item Select instance to explain
				\item Generate perturbed samples around it
				\item Get predictions from complex model
				\item Fit interpretable model on perturbed data
				\item Use interpretable model to explain prediction
			\end{enumerate}
			\item \textbf{Key features:}
			\begin{itemize}
				\item \textbf{Local:} Only accurate near the instance
				\item \textbf{Model-agnostic:} Works with any model
				\item \textbf{Interpretable:} Uses simple model
				\item \textbf{Sparse:} Uses few features
			\end{itemize}
			\item \textbf{Example:}
			\begin{itemize}
				\item Complex DNN predicted loan denial
				\item LIME says: low income and high debt caused denial
				\item Explanation is local to that case
			\end{itemize}
		\end{itemize}
		
		\vspace{0.05cm}
		\textcolor{myred}{\textbf{WHY OTHER OPTIONS ARE WRONG:}}
		\begin{itemize}
			\item \textcolor{myred}{B:} LIME is for explanation, not accuracy
			\item \textcolor{myred}{C:} LIME is not regularization
			\item \textcolor{myred}{D:} LIME is not feature selection (though it identifies important features locally)
		\end{itemize}
		
		\textcolor{myred}{\textbf{EXAM TRAP:}} LIME provides local explanations for individual predictions!
	\end{frame}
	
	\begin{frame}{Q100: Model Estimation Summary - Tutorial}
		\fontsize{6.5}{7.5}\selectfont
		\textbf{Which provides the most comprehensive summary of model estimation?}
		
		\vspace{0.05cm}
		\textcolor{myblue}{\textbf{ANSWER: B}} - Model estimation involves choosing between OLS, NLS, and MLE, optimizing using gradient descent, and evaluating with cross-validation.
		
		\vspace{0.05cm}
		\textbf{DETAILED EXPLANATION:}
		\begin{itemize}
			\item \textbf{Complete picture of model estimation:}
			\begin{enumerate}
				\item \textbf{Estimation method selection:}
				\begin{itemize}
					\item OLS: linear in parameters
					\item NLS: nonlinear in parameters
					\item MLE: maximum likelihood
				\end{itemize}
				\item \textbf{Optimization:}
				\begin{itemize}
					\item Gradient descent (and variants)
					\item Backpropagation (neural networks)
					\item Numerical methods for NLS
				\end{itemize}
				\item \textbf{Evaluation:}
				\begin{itemize}
					\item Cross-validation
					\item Bootstrap
					\item Performance metrics
				\end{itemize}
				\item \textbf{Regularization:}
				\begin{itemize}
					\item L1/L2 regularization
					\item Dropout
					\item Early stopping
				\end{itemize}
			\end{enumerate}
			\item \textbf{Key insight:} All components work together
			\begin{itemize}
				\item Each step is important
				\item Choices affect final model
				\item Tradeoffs at every step
			\end{itemize}
		\end{itemize}
		
		\vspace{0.05cm}
		\textcolor{myred}{\textbf{WHY OTHER OPTIONS ARE WRONG:}}
		\begin{itemize}
			\item \textcolor{myred}{A:} Only OLS is incomplete
			\item \textcolor{myred}{C:} Only MLE is incomplete
			\item \textcolor{myred}{D:} Only regularization is incomplete
		\end{itemize}
		
		\textcolor{myred}{\textbf{EXAM TRAP:}} Model estimation is comprehensive - includes method choice, optimization, and evaluation!
	\end{frame}
	
	% ============================================================
	% SUMMARY SLIDE
	% ============================================================
	
	\begin{frame}{Summary: Complete Review}
		\fontsize{6.5}{7.5}\selectfont
		\textbf{\textcolor{myblue}{Core Estimation Methods:}}
		\begin{itemize}
			\item \textcolor{myblue}{OLS:} Minimizes SSE; closed-form for linear regression; BLUE under Gauss-Markov
			\item \textcolor{myblue}{MLE:} Maximizes likelihood; equivalent to OLS under normal errors
			\item \textcolor{myblue}{NLS:} For nonlinear-in-parameters models; requires numerical optimization
		\end{itemize}
		
		\vspace{0.1cm}
		\textbf{\textcolor{myblue}{Optimization:}}
		\begin{itemize}
			\item \textcolor{myblue}{Gradient Descent:} Iterative; uses negative gradient; learning rate controls step size
			\item \textcolor{myblue}{Backpropagation:} Chain rule for neural networks; requires storing activations
			\item \textcolor{myblue}{Adam/Momentum:} Adaptive methods for faster convergence
		\end{itemize}
		
		\vspace{0.1cm}
		\textbf{\textcolor{myblue}{Model Selection:}}
		\begin{itemize}
			\item \textcolor{myblue}{Bias-Variance Tradeoff:} Simplicity vs complexity; optimal at minimum total error
			\item \textcolor{myblue}{Regularization:} L1 (Lasso - feature selection), L2 (Ridge - shrinkage)
			\item \textcolor{myblue}{Cross-Validation:} k-fold (k=10 common); evaluates generalization
		\end{itemize}
		
		\vspace{0.1cm}
		\textbf{\textcolor{myblue}{Interpretability:}}
		\begin{itemize}
			\item \textcolor{myblue}{White-box models:} Decision trees, linear regression (interpretable)
			\item \textcolor{myblue}{Black-box models:} Deep neural networks (accurate but not interpretable)
			\item \textcolor{myblue}{Explainability:} SHAP values, LIME for understanding predictions
		\end{itemize}
		
		\vspace{0.1cm}
		\textcolor{myred}{\textbf{Exam Success:}} Understand the tradeoffs and know the definitions!
	\end{frame}
	
\end{document}